% ============================================================================
% Statistical validation section for the multi-year CafeF--FinABSA study.
% This is a section fragment: include it from the main document with
% % ============================================================================
% Statistical validation section for the multi-year CafeF--FinABSA study.
% This is a section fragment: include it from the main document with
% % ============================================================================
% Statistical validation section for the multi-year CafeF--FinABSA study.
% This is a section fragment: include it from the main document with
% % ============================================================================
% Statistical validation section for the multi-year CafeF--FinABSA study.
% This is a section fragment: include it from the main document with
% \input{statistical_validation_multiyear.tex}.
%
% Required packages in the main preamble:
% amsmath, amssymb, booktabs, graphicx, longtable, array, hyperref
%
% Figures are expected under images/.
% All numerical values below are from the completed October 2019--2022 runs.
% ============================================================================

\section{Empirical Validation and Statistical Testing}
\label{sec:empirical-validation}

This section evaluates whether target-specific sentiment predictions generated by the FinABSA pipeline are associated with subsequent abnormal stock returns in Vietnamese financial news. The analysis extends the original October 2022 design to four comparable October samples: October 2019, October 2020, October 2021, and October 2022. October 2019 is retained as a pre-pandemic reference period, while 2020--2022 capture different stages of the COVID-19 and post-COVID market environment. The same calendar month is used in every year to reduce month-of-year confounding. Each period is analysed separately using its own price history from May through December of the same year.

The empirical objective is deliberately narrower than a causal claim. The tests ask whether predicted sentiment is statistically associated with market-adjusted price movements after entity linking, time alignment, and aggregation to the ticker--day level. They do not establish that a news article caused a return, that the model generates a profitable trading strategy, or that the observed association generalises to all Vietnamese financial news.

\subsection{Research Questions and Hypotheses}
\label{subsec:research-questions-statistics}

The downstream validation addresses three questions. First, are the model outputs complete and correctly aligned with the target-masked article--ticker inputs? Second, do abnormal returns differ from zero around the first trading session after publication? Third, do the results remain stable across sentiment groups, control samples, fixed-effects specifications, and lag definitions?

The primary null hypotheses are:
\begin{itemize}
    \item $H_{0,1}$: the mean cumulative abnormal return is zero within each event window;
    \item $H_{0,2}$: abnormal returns do not differ between eligible ticker--days with an identified news event and eligible ticker--days without one;
    \item $H_{0,3}$: abnormal returns do not differ systematically across positive, neutral, and negative sentiment groups;
    \item $H_{0,4}$: the conditional association between sentiment and abnormal return is zero after ticker and date fixed effects;
    \item $H_{0,5}$: any observed association does not persist in a stable and monotonic pattern across alternative lags.
\end{itemize}

All hypothesis tests are two-sided and use a nominal significance level of $\alpha=0.05$. Because many period--window--group combinations are examined, individual $p$-values are treated as exploratory unless the pattern is replicated across specifications and survives an appropriate multiple-comparison correction.

\subsection{Data Construction and Multi-year Coverage}
\label{subsec:multi-year-data}

\paragraph{News corpus.}
The crawler collected historical CafeF articles from the October sitemap partitions for each year. The resulting corpus contains 5,133 article records. A stable article identifier is derived from the URL, and each record retains the title, publication timestamp, URL, parsing status, and related metadata. The crawl status is kept explicitly so that an unavailable or malformed page is not silently interpreted as an empty financial article.

A headline can mention more than one company. Since FinABSA predicts sentiment for a specified target rather than discovering ticker symbols itself, the inference unit is an article--ticker pair. Entity linking produces a broad candidate file and a strict high-confidence file. The broad file retains all automatically detected candidates, while the strict file retains candidates with entity-linking confidence at least 0.95. This distinction is a measurement-quality choice rather than a model change.

\begin{table}[htbp]
\centering
\small
\caption{Multi-year CafeF corpus and model-input coverage.}
\label{tab:multi-year-news-coverage}
\begin{tabular}{lrrrrrr}
\toprule
\textbf{Period} & \textbf{Articles} & \textbf{Crawl OK} & \textbf{Full pairs} & \textbf{Strict pairs} & \textbf{Strict + price} & \textbf{Price errors} \\
\midrule
October 2019 & 1,040 & 1,039 & 583 & 154 & 141 & 11 \\
October 2020 & 1,115 & 1,115 & 657 & 182 & 168 & 10 \\
October 2021 & 1,506 & 1,491 & 991 & 297 & 281 & 11 \\
October 2022 & 1,472 & 1,457 & 509 & 206 & 206 & 0 \\
\midrule
\textbf{Total} & \textbf{5,133} & -- & \textbf{2,740} & \textbf{839} & \textbf{796} & \textbf{32} \\
\bottomrule
\end{tabular}
\end{table}

The price errors are not replaced by synthetic observations. They correspond to candidate strings that could not be retrieved as valid securities or to failed endpoint requests. Similarly, the 43 strict input rows without a usable price window are preserved in the inference file but excluded from financial calculations that require a matched price series. This separation prevents missing market observations from being treated as zero returns.

\paragraph{Target masking and prediction contract.}
For each article--ticker pair, the target surface form is replaced with \texttt{Target} and other detected entities are replaced with \texttt{Other}. The model therefore receives a target-masked headline and produces one of three canonical labels: positive, neutral, or negative. The model does not perform named entity recognition or ticker prediction; entity linking is a separate upstream step.

\begin{table}[htbp]
\centering
\small
\caption{Prediction validation by October period.}
\label{tab:multi-year-prediction-validation}
\begin{tabular}{lrrrrrrrr}
\toprule
\textbf{Period} & \textbf{Input} & \textbf{Pred.} & \textbf{Missing} & \textbf{Unexpected} & \textbf{Dup.} & \textbf{Unknown} & \textbf{Neutral} & \textbf{Pos./Neg.} \\
\midrule
October 2019 & 154 & 154 & 0 & 0 & 0 & 0 & 144 & 7 / 3 \\
October 2020 & 182 & 182 & 0 & 0 & 0 & 0 & 167 & 13 / 2 \\
October 2021 & 297 & 297 & 0 & 0 & 0 & 0 & 270 & 16 / 11 \\
October 2022 & 206 & 206 & 0 & 0 & 0 & 0 & 178 & 18 / 10 \\
\midrule
\textbf{Total} & \textbf{839} & \textbf{839} & \textbf{0} & \textbf{0} & \textbf{0} & \textbf{0} & \textbf{759} & \textbf{54 / 26} \\
\bottomrule
\end{tabular}
\end{table}

The prediction file is aligned to the input contract by the stable \texttt{sample\_id}, not by row order. All 839 strict inputs have exactly one matching prediction. There are no missing predictions, unexpected prediction IDs, duplicate IDs, or unknown labels after canonicalisation. This validates the data plumbing, but it does not validate semantic accuracy: without independently annotated gold labels, the prediction distribution cannot be converted into accuracy, macro-F1, MCC, or Cohen's kappa.

\begin{figure}[htbp]
\centering
\includegraphics[width=0.78\textwidth]{images/multiyear_label_distribution.png}
\caption{Composition of the predicted sentiment labels by October period. Neutral predictions dominate every period, while positive and negative groups are comparatively small. The imbalance reduces the precision of group-level comparisons and motivates reporting group sizes together with every estimate.}
\label{fig:multi-year-label-distribution}
\end{figure}

\paragraph{Market data.}
Daily OHLCV data were retrieved through \texttt{vnstock} for the identified securities and the VNINDEX and VN30 reference series. For every October period, the stored price window begins in May of the same year and ends in December. This provides pre-event observations for the market model and post-event observations for lag calculations without importing prices from a different calendar regime.

\begin{table}[htbp]
\centering
\small
\caption{Market-price coverage by October period.}
\label{tab:multi-year-price-coverage}
\begin{tabular}{lrrrrr}
\toprule
\textbf{Period} & \textbf{Price rows} & \textbf{Symbols} & \textbf{Price start} & \textbf{Price end} & \textbf{Strict rows with price} \\
\midrule
October 2019 & 8,320 & 49 & 2019-05-02 & 2019-12-31 & 141 \\
October 2020 & 9,449 & 57 & 2020-05-04 & 2020-12-31 & 168 \\
October 2021 & 11,935 & 73 & 2021-05-04 & 2021-12-31 & 281 \\
October 2022 & 9,182 & 55 & 2022-05-04 & 2022-12-30 & 206 \\
\bottomrule
\end{tabular}
\end{table}

The number of symbols differs by period because the news corpus and entity-linking output differ across years. The price data are therefore comparable in construction, but not identical in cross-sectional composition. This point is important when comparing the magnitude of estimates across periods.

\subsection{Time Alignment and Return Construction}
\label{subsec:time-alignment-multi-year}

The data contain daily prices but not the exact time at which an article was disseminated to market participants. The implemented alignment rule maps each article's normalized publication date to the first trading date strictly after that calendar date. Thus, an article published on a trading day is evaluated against the next available trading session; an article published on a weekend or holiday is also mapped to the next trading session. No fixed 15:00 cutoff is imposed, because the daily data do not identify the intraday time at which the information was incorporated into the price.

For security $i$ on trading date $t$, the continuously compounded percentage return is computed as
\begin{equation}
 r_{i,t}=100\times\left[\log(P_{i,t})-\log(P_{i,t-1})\right],
\end{equation}
where $P_{i,t}$ is the recorded closing price. The market return is computed using the same transformation for VNINDEX. The analysis uses the downloaded price series as stored; any interpretation of return magnitude is conditional on the source's adjustment and corporate-action coverage.

The sentiment score is a deterministic label encoding:
\begin{equation}
 S_{i,t}=\begin{cases}
 +1,&\text{positive prediction},\\
 0,&\text{neutral prediction},\\
 -1,&\text{negative prediction}.
 \end{cases}
\end{equation}

Several article--ticker rows can refer to the same ticker on the same date. Because the market supplies one return observation per ticker--day, the financial unit is formed by aggregation before inference. The panel records mean sentiment, positive/neutral/negative shares, the number of unique articles, and the number of contributing article--ticker rows. This avoids artificially multiplying the effective sample size by repeating one return for every article.

\subsection{Event-Study Methodology}
\label{subsec:event-study-multi-year}

Following the standard market-model event-study framework of Brown and Warner~\cite{brownwarner1985}, the expected return for each security is estimated from observations before 1 October of the relevant year:
\begin{equation}
 R_{i,t}=\alpha_i+\beta_iR_{m,t}+\varepsilon_{i,t},
\end{equation}
where $R_{m,t}$ is the VNINDEX return. The abnormal return is
\begin{equation}
 AR_{i,t}=R_{i,t}-\left(\widehat{\alpha}_i+\widehat{\beta}_iR_{m,t}\right).
\end{equation}
A security-specific market model is fitted when at least 20 pre-event observations are available; otherwise the documented fallback parameters are retained. The final period-specific runs contain 49, 57, 73, and 55 stored price symbols, respectively, with 54 fitted market-model records in the October 2022 run and corresponding period-specific records in the other runs.

For an event aligned to trading date $t$, cumulative abnormal return is defined as
\begin{equation}
 CAR_i(T_1,T_2)=\sum_{\tau=T_1}^{T_2}AR_{i,t+\tau}.
\end{equation}
Three windows are examined: $[0,0]$ for the event session, $[0,+1]$ for the event session and the next session, and $[0,+3]$ for the event session and the next three sessions. For every period and window, the pipeline reports the mean and median CAR, a two-sided one-sample $t$-test against zero, a two-sided sign-flip permutation test with 5,000 replications, and a percentile bootstrap 95\% confidence interval with 5,000 replications. Fixed random seeds make the resampling results reproducible.

\begin{figure}[htbp]
\centering
\includegraphics[width=0.84\textwidth]{images/multiyear_car_by_period.png}
\caption{Aggregate CAR by October period and event window. Points show mean CAR and whiskers show percentile bootstrap 95\% confidence intervals. The figure is intended for cross-period comparison; it does not apply a multiple-comparison correction and does not establish causality.}
\label{fig:multi-year-car}
\end{figure}

\begin{longtable}{llrrrrrr}
\caption{Aggregate event-study CAR tests by period and event window.}\label{tab:multi-year-car-tests}\\
\toprule
\textbf{Period} & \textbf{Window} & \textbf{N} & \textbf{Mean} & \textbf{Median} & \textbf{$p_t$} & \textbf{$p_{perm}$} & \textbf{Bootstrap 95\% CI} \\
\midrule
\endfirsthead
\toprule
\textbf{Period} & \textbf{Window} & \textbf{N} & \textbf{Mean} & \textbf{Median} & \textbf{$p_t$} & \textbf{$p_{perm}$} & \textbf{Bootstrap 95\% CI} \\
\midrule
\endhead
\midrule
\multicolumn{8}{r}{Continued on next page}\\
\endfoot
\bottomrule
\endlastfoot
2019-10 & $[0,0]$ & 120 & -0.1323 & -0.1570 & 0.5420 & 0.5491 & $[-0.5550,0.2916]$ \\
2019-10 & $[0,+1]$ & 120 & 0.0573 & -0.1285 & 0.8833 & 0.8830 & $[-0.6841,0.8108]$ \\
2019-10 & $[0,+3]$ & 120 & 0.0020 & -0.2812 & 0.9970 & 0.9974 & $[-1.0491,1.0923]$ \\
2020-10 & $[0,0]$ & 134 & 0.0522 & -0.2528 & 0.8031 & 0.8072 & $[-0.3645,0.4595]$ \\
2020-10 & $[0,+1]$ & 134 & -0.2975 & -0.6585 & 0.3424 & 0.3363 & $[-0.8991,0.3050]$ \\
2020-10 & $[0,+3]$ & 133 & -0.9261 & -1.3113 & 0.0438 & 0.0440 & $[-1.8254,-0.0657]$ \\
2021-10 & $[0,0]$ & 212 & -0.2797 & -0.4321 & 0.1070 & 0.1096 & $[-0.6117,0.0587]$ \\
2021-10 & $[0,+1]$ & 212 & -0.2584 & -0.5174 & 0.2920 & 0.2953 & $[-0.7377,0.2377]$ \\
2021-10 & $[0,+3]$ & 212 & -1.0149 & -1.0412 & 0.0101 & 0.0088 & $[-1.7967,-0.2662]$ \\
2022-10 & $[0,0]$ & 150 & 0.1363 & 0.0008 & 0.6023 & 0.6091 & $[-0.3557,0.6421]$ \\
2022-10 & $[0,+1]$ & 150 & -0.0220 & 0.1798 & 0.9501 & 0.9562 & $[-0.7298,0.6360]$ \\
2022-10 & $[0,+3]$ & 150 & -0.2889 & -0.7977 & 0.5614 & 0.5615 & $[-1.2725,0.6793]$ \\
\end{longtable}

\paragraph{Event-study results.}
The aggregate CAR results do not support a single stable sentiment-linked price response across the four October periods. The same-day window is statistically indistinguishable from zero in every period: all $p$-values exceed 0.10 and all bootstrap intervals include zero. In the three-session window, 2020-10 and 2021-10 produce negative mean CAR estimates with nominal $p$-values below 0.05. These two results are noteworthy but should not be treated as decisive evidence: they are two findings among many period--window tests, they are not adjusted for multiple comparisons, and the negative three-session CAR is not itself evidence that sentiment caused the decline.

The 2019 and 2022 estimates remain close to zero or statistically imprecise across all windows. This cross-period heterogeneity is more consistent with a period-dependent market association than with a universal FinABSA effect. It may reflect differences in market regime, news composition, entity coverage, macroeconomic shocks, or the composition of the available price universe.

\subsection{News-Day Versus No-News-Day Control}
\label{subsec:news-control-multi-year}

The event study is complemented by a control comparison. For each period, the control sample contains eligible ticker--days belonging to tickers that appear in the event sample, restricted to October of that year. A news ticker--day is a ticker--day with at least one aligned article event; a no-news ticker--day is an eligible ticker--day without an identified event. The estimand is
\begin{equation}
 \Delta=\overline{AR}_{news}-\overline{AR}_{no\text{-}news}.
\end{equation}

The two distributions are compared using a Welch unequal-variance test, a pooled-label permutation test with 5,000 permutations, and a percentile bootstrap interval for the difference in means.

\begin{table}[htbp]
\centering
\scriptsize
\caption{News-day versus eligible no-news-day abnormal returns.}
\label{tab:news-control-multi-year}
\begin{tabular}{lrrrrrrrr}
\toprule
\textbf{Period} & \textbf{$N_{news}$} & \textbf{$N_{0}$} & \textbf{News mean} & \textbf{No-news mean} & \textbf{$\Delta$} & \textbf{Welch $p$} & \textbf{Perm. $p$} & \textbf{Bootstrap 95\% CI} \\
\midrule
2019-10 & 120 & 941 & -0.1323 & 0.0064 & -0.1386 & 0.5383 & 0.4585 & $[-0.5937,0.3095]$ \\
2020-10 & 134 & 1,027 & 0.0522 & -0.2161 & 0.2683 & 0.2240 & 0.1828 & $[-0.1601,0.7062]$ \\
2021-10 & 212 & 1,217 & -0.2797 & -0.1039 & -0.1758 & 0.3553 & 0.3771 & $[-0.5493,0.1980]$ \\
2022-10 & 150 & 936 & 0.1363 & -0.3559 & 0.4921 & 0.0804 & 0.0684 & $[-0.0645,1.0442]$ \\
\bottomrule
\end{tabular}
\end{table}

\begin{figure}[htbp]
\centering
\includegraphics[width=0.78\textwidth]{images/multiyear_news_day_control.png}
\caption{News-day minus eligible no-news-day abnormal return by October period. Whiskers show percentile bootstrap 95\% confidence intervals; each interval includes zero.}
\label{fig:multi-year-news-control}
\end{figure}

The control comparison does not reject equality at the 5\% level in any period. The largest positive point estimate occurs in October 2022, where news days have 0.4921 percentage points higher abnormal returns than no-news days, but the Welch and permutation $p$-values are 0.0804 and 0.0684, respectively, and the bootstrap interval crosses zero. The result is therefore suggestive but not conventionally significant. The comparison remains observational because news selection is not random: firms that receive coverage may systematically differ from firms that do not.

\subsection{Sentiment-Group Contrasts}
\label{subsec:sentiment-groups-multi-year}

At the ticker--day level, the analysis compares three pre-specified contrasts: positive minus neutral, negative minus neutral, and positive minus negative. Each contrast is evaluated with a Welch test, a pooled permutation test, and a bootstrap interval. Differences are defined as group A minus group B.

\begin{longtable}{llrrrrrr}
\caption{Sentiment-group abnormal-return contrasts by period.}\label{tab:sentiment-contrasts-multi-year}\\
\toprule
\textbf{Period} & \textbf{Contrast} & \textbf{$N_A$} & \textbf{$N_B$} & \textbf{Difference} & \textbf{Welch $p$} & \textbf{Perm. $p$} & \textbf{Bootstrap 95\% CI} \\
\midrule
\endfirsthead
\toprule
\textbf{Period} & \textbf{Contrast} & \textbf{$N_A$} & \textbf{$N_B$} & \textbf{Difference} & \textbf{Welch $p$} & \textbf{Perm. $p$} & \textbf{Bootstrap 95\% CI} \\
\midrule
\endhead
\midrule
\multicolumn{8}{r}{Continued on next page}\\
\endfoot
\bottomrule
\endlastfoot
2019-10 & Positive $-$ Neutral & 6 & 113 & -1.5118 & 0.1196 & 0.1252 & $[-3.2365,-0.3238]$ \\
2019-10 & Negative $-$ Neutral & 1 & 113 & -- & -- & -- & Not estimable \\
2019-10 & Positive $-$ Negative & 6 & 1 & -- & -- & -- & Not estimable \\
2020-10 & Positive $-$ Neutral & 10 & 123 & -0.8544 & 0.3380 & 0.2813 & $[-2.1566,0.9352]$ \\
2020-10 & Negative $-$ Neutral & 1 & 123 & -- & -- & -- & Not estimable \\
2020-10 & Positive $-$ Negative & 10 & 1 & -- & -- & -- & Not estimable \\
2021-10 & Positive $-$ Neutral & 12 & 191 & 1.5693 & 0.0267 & 0.0388 & $[0.5210,2.8516]$ \\
2021-10 & Negative $-$ Neutral & 9 & 191 & 0.0915 & 0.9093 & 0.9112 & $[-1.0946,1.7607]$ \\
2021-10 & Positive $-$ Negative & 12 & 9 & 1.4779 & 0.1460 & 0.1372 & $[-0.4072,3.2104]$ \\
2022-10 & Positive $-$ Neutral & 14 & 126 & 0.6990 & 0.4518 & 0.4271 & $[-0.9308,2.4970]$ \\
2022-10 & Negative $-$ Neutral & 10 & 126 & -0.4688 & 0.6712 & 0.6601 & $[-2.5243,1.4876]$ \\
2022-10 & Positive $-$ Negative & 14 & 10 & 1.1678 & 0.3962 & 0.3971 & $[-1.2938,3.6937]$ \\
\end{longtable}

The most notable group contrast appears in October 2021: the positive group has a 1.5693 percentage-point higher mean abnormal return than the neutral group, with Welch $p=0.0267$, permutation $p=0.0388$, and a bootstrap interval that excludes zero. This result should be described as a period-specific exploratory finding rather than a general law. The positive group contains only 12 ticker--day observations, the negative group contains nine, the analysis involves multiple comparisons, and the effect is not replicated in the corresponding positive--neutral comparisons for 2019, 2020, or 2022.

The 2019 positive--neutral point estimate is negative, whereas the 2020 and 2022 differences are also statistically imprecise. The failure of sign consistency across years weakens the claim that predicted sentiment produces a stable cross-sectional ranking of abnormal returns.

\begin{figure}[htbp]
\centering
\includegraphics[width=0.78\textwidth]{images/multiyear_positive_neutral_contrast.png}
\caption{Positive-minus-neutral abnormal-return contrast by October period. Whiskers show bootstrap 95\% confidence intervals. The positive 2021 estimate is not sufficient on its own to establish a stable sentiment effect across years.}
\label{fig:multi-year-positive-neutral}
\end{figure}

\subsection{Two-Way Fixed-Effects Panel Regression}
\label{subsec:panel-regression-multi-year}

To control for persistent differences between securities and common shocks across dates, the following two-way fixed-effects model is estimated separately for each October period:
\begin{equation}
 AR_{i,\tau(i,t)}=\alpha_i+\delta_{\tau(i,t)}+\beta_1S_{i,t}
 +\beta_2\mathrm{NegativeShare}_{i,t}
 +\beta_3\mathrm{PositiveShare}_{i,t}
 +\gamma\log(1+\mathrm{ArticleCount}_{i,t})+\varepsilon_{i,t},
\end{equation}
where $\alpha_i$ denotes ticker fixed effects, $\delta_{\tau(i,t)}$ denotes date fixed effects, and $\tau(i,t)$ is the first trading session strictly after publication. The dependent variable is market-model abnormal return on the aligned evaluation date.

\begin{table}[htbp]
\centering
\scriptsize
\caption{Key coefficients from period-specific two-way fixed-effects regressions. The full coefficient tables, including ticker and date fixed-effect indicators, are provided in the accompanying CSV output.}
\label{tab:panel-key-coefficients-multi-year}
\begin{tabular}{llrrrr}
\toprule
\textbf{Period} & \textbf{Variable} & \textbf{Coefficient} & \textbf{Std. error} & \textbf{$t$-statistic} & \textbf{$p$-value} \\
\midrule
2019-10 & Sentiment score & -0.4450 & 0.8508 & -0.5230 & 0.6009 \\
2019-10 & Negative share & 1.5084 & 1.8412 & 0.8192 & 0.4127 \\
2019-10 & Positive share & 1.0633 & 1.2562 & 0.8468 & 0.3973 \\
2019-10 & Log article count & -1.5405 & 3.3930 & -0.4540 & 0.6498 \\
2020-10 & Sentiment score & 2.5056 & 0.8042 & 3.1156 & 0.0018 \\
2020-10 & Negative share & -4.9332 & 1.8033 & -2.7356 & 0.0062 \\
2020-10 & Positive share & -2.4276 & 1.2770 & -1.9010 & 0.0573 \\
2020-10 & Log article count & 2.3897 & 1.9986 & 1.1957 & 0.2318 \\
2021-10 & Sentiment score & -0.9696 & 2.4681 & -0.3927 & 0.6944 \\
2021-10 & Negative share & 0.1799 & 2.5174 & 0.0715 & 0.9430 \\
2021-10 & Positive share & 0.4345 & 2.6128 & 0.1663 & 0.8679 \\
2021-10 & Log article count & 0.5867 & 1.0124 & 0.5795 & 0.5623 \\
2022-10 & Sentiment score & 6.5062 & 5.3763 & 1.2100 & 0.2262 \\
2022-10 & Negative share & 5.9827 & 5.2919 & 1.1305 & 0.2582 \\
2022-10 & Positive share & -6.1476 & 5.5086 & -1.1160 & 0.2644 \\
2022-10 & Log article count & 0.3079 & 1.3849 & 0.2224 & 0.8241 \\
\bottomrule
\end{tabular}
\end{table}

The sentiment-score coefficient is positive and statistically significant only in the 2020 period. The coefficient is negative and insignificant in 2019 and 2021, and positive but imprecise in 2022. This sign instability is central: a robust sentiment--return association would be expected to exhibit a more consistent direction under comparable definitions. The 2020 coefficient therefore constitutes a period-specific association, not evidence that the model has a stable return-prediction relationship.

The coefficient on negative share is also significant only in 2020 and is negative, while the positive-share coefficient is not significant at the 5\% level. Because sentiment shares are compositional and mechanically related to the omitted neutral share, their coefficients should not be interpreted as independent causal marginal effects. The relatively high within-period $R^2$ values are not evidence of strong predictive power: ticker and date fixed effects can absorb substantial variation, while the sentiment covariates remain sensitive to sample composition.

The intended covariance estimator was HC3. In sparse two-way fixed-effects designs, HC3 can produce infinite standard errors when leverage reaches a singular boundary. The implementation therefore falls back to HC1 when necessary and records this fallback in the period-specific metadata. These panel coefficients should consequently be treated as exploratory conditional associations.

\begin{figure}[htbp]
\centering
\includegraphics[width=0.78\textwidth]{images/multiyear_panel_sentiment_coefficient.png}
\caption{Period-specific sentiment-score coefficients from the two-way fixed-effects models. Whiskers show coefficient plus or minus 1.96 standard errors. The 2020 coefficient is positive and relatively precise, but the sign is not replicated across the other periods.}
\label{fig:multi-year-panel-coefficient}
\end{figure}

\subsection{Lagged Robustness and Sensitivity Analysis}
\label{subsec:lag-robustness-multi-year}

Lag analysis evaluates abnormal returns at 0, 1, 2, 3, 5, and 10 trading sessions after the aligned event date. These are sensitivity checks rather than independent causal tests. The sample size decreases at longer lags because the relevant security price must remain available after the event.

\begin{longtable}{llrrrrrr}
\caption{Lagged abnormal-return robustness by period.}\label{tab:lag-robustness-multi-year}\\
\toprule
\textbf{Period} & \textbf{Lag} & \textbf{N} & \textbf{Overall AR} & \textbf{Negative} & \textbf{Neutral} & \textbf{Positive} \\
\midrule
\endfirsthead
\toprule
\textbf{Period} & \textbf{Lag} & \textbf{N} & \textbf{Overall AR} & \textbf{Negative} & \textbf{Neutral} & \textbf{Positive} \\
\midrule
\endhead
\midrule
\multicolumn{7}{r}{Continued on next page}\\
\endfoot
\bottomrule
\endlastfoot
2019-10 & 0 & 120 & -0.1323 & -1.5766 & -0.0648 & 0.8940 \\
2019-10 & 1 & 146 & 0.3194 & 0.0588 & 0.2772 & 1.2785 \\
2019-10 & 2 & 141 & -0.0968 & 0.0331 & -0.0974 & -0.1286 \\
2019-10 & 3 & 134 & 0.0067 & -0.4507 & 0.0283 & -0.3248 \\
2019-10 & 5 & 108 & 0.1154 & -0.6470 & 0.1166 & 0.2232 \\
2019-10 & 10 & 74 & -0.1635 & 0.1766 & -0.1547 & -0.6273 \\
2020-10 & 0 & 134 & 0.2261 & -1.9519 & 0.3085 & -0.4498 \\
2020-10 & 1 & 164 & -0.1520 & 0.5470 & -0.1975 & 0.2915 \\
2020-10 & 2 & 159 & -0.3772 & 1.3165 & -0.3780 & -0.7051 \\
2020-10 & 3 & 152 & 0.0687 & 0.9315 & 0.1112 & -0.6345 \\
2020-10 & 5 & 144 & -0.0997 & 0.4069 & -0.1472 & 0.4869 \\
2020-10 & 10 & 100 & -0.4924 & 0.3574 & -0.5596 & 1.5471 \\
2021-10 & 0 & 212 & -0.0285 & -0.1815 & -0.0939 & 1.0708 \\
2021-10 & 1 & 267 & 0.2065 & 0.6397 & 0.2368 & -0.5095 \\
2021-10 & 2 & 255 & -0.3156 & -0.4267 & -0.2764 & -0.8200 \\
2021-10 & 3 & 244 & 0.0773 & -0.3897 & 0.1940 & -1.2377 \\
2021-10 & 5 & 219 & 0.1398 & -0.1027 & 0.1966 & -0.4387 \\
2021-10 & 10 & 147 & -0.0471 & 1.2318 & -0.1101 & -0.1871 \\
2022-10 & 0 & 150 & -0.0936 & -0.7884 & -0.1653 & 0.8940 \\
2022-10 & 1 & 174 & -0.1797 & 0.3368 & 0.0713 & -2.7331 \\
2022-10 & 2 & 163 & -0.2007 & -1.2344 & -0.0991 & -0.4067 \\
2022-10 & 3 & 153 & -0.5972 & -2.6554 & -0.4837 & -0.1760 \\
2022-10 & 5 & 134 & -0.7243 & 0.4479 & -0.9416 & 0.5251 \\
2022-10 & 10 & 93 & -0.5183 & -0.4422 & -0.2779 & -4.1530 \\
\end{longtable}

\begin{figure}[htbp]
\centering
\includegraphics[width=0.84\textwidth]{images/multiyear_lag_robustness.png}
\caption{Lagged sensitivity of mean abnormal returns across the four October periods. The changing signs and magnitudes indicate that the association is sensitive to the evaluation horizon rather than a stable delayed response.}
\label{fig:multi-year-lag-robustness}
\end{figure}

The lag results do not display a common monotonic decay. In 2019, the overall estimate changes from negative at lag 0 to positive at lag 1 and back toward zero thereafter. In 2020, the estimates become more negative at lags 2 and 10, while in 2021 the sign alternates across horizons. In 2022, the overall mean is increasingly negative at lags 3 and 5. Sentiment-specific estimates are even less stable, particularly for the small positive and negative groups. The appropriate conclusion is that the apparent price association depends on the timing convention and the available evaluation horizon.

\subsection{Additional Diagnostic Evidence}
\label{subsec:additional-diagnostics-multi-year}

The output package also contains daily ticker--day sentiment and abnormal-return tables, market-model coverage, event observations, and panel-regression data. These artifacts support auditability and allow the reader to trace each reported result back to an aligned ticker--day observation. The following diagnostic figures are descriptive rather than inferential.

\begin{figure}[htbp]
\centering
\includegraphics[width=0.82\textwidth]{images/market_model_coverage.png}
\caption{Market-model coverage diagnostic from the completed October 2022 run. It shows the number of pre-event observations available for individual symbols and illustrates why coverage thresholds and fallback handling are necessary.}
\label{fig:market-model-coverage-multi-year}
\end{figure}

\begin{figure}[htbp]
\centering
\includegraphics[width=0.90\textwidth]{images/sentiment_vs_abnormal_return.png}
\caption{Descriptive comparison between daily mean sentiment and mean abnormal return in the ticker--day panel. The figure is not a causal identification plot and does not control for all market-wide shocks.}
\label{fig:sentiment-vs-ar-multi-year}
\end{figure}

\subsection{Multiple Comparisons and Statistical Interpretation}
\label{subsec:multiple-comparisons}

The multi-year design produces several families of related tests: four periods, three event windows, multiple sentiment groups, a news-day control, panel coefficients, and six lag horizons. Consequently, at least some nominally small $p$-values can occur by chance when many hypotheses are examined. The results are therefore interpreted using three criteria: the size and direction of the estimate, whether the confidence interval excludes zero, and whether the finding is replicated under related specifications or periods.

Under this framework, the 2020 and 2021 three-day CAR findings and the 2021 positive--neutral contrast are interesting signals for further study, but not definitive discoveries. They should be reported with their exact sample sizes, intervals, and unadjusted $p$-values. A confirmatory study should pre-specify a smaller number of primary hypotheses and apply a family-wise or false-discovery-rate correction, such as Holm or Benjamini--Hochberg, depending on the inferential objective.

\subsection{Analyses Not Validly Estimated in the Current Run}
\label{subsec:deferred-analyses-multi-year}

Several analyses appeared in the initial experimental plan but are not assigned numerical claims here because the required artifacts are absent. First, intrinsic CafeF accuracy, macro-F1, MCC, and Cohen's kappa require an independently human-annotated gold-label sample. The model's own predictions cannot serve as their ground truth. Second, probability calibration metrics such as Expected Calibration Error require valid per-class probabilities or logits; the current prediction files contain hard labels and empty probability fields. Third, out-of-time return forecasting requires a pre-specified chronological training, validation, and test design with a fitted return-prediction model; the current pipeline is an event-study association analysis, not a trained forecaster. Fourth, a broad-linking robustness result requires a complete independent inference and financial analysis run on the 2,740-row broad input set; strict and full input files alone are not sufficient.

This separation is methodologically important. It prevents unsupported accuracy, calibration, forecasting, and robustness numbers from being presented as results. The completed outputs can be extended later without changing the core data contract.

\subsection{Limitations}
\label{subsec:limitations-statistical}

\paragraph{Limited temporal coverage and regime dependence.}
Although the study now covers four October periods, it still samples only one month per year. October 2019--2022 may contain unusual news intensity, market stress, or sector composition. The estimates should not be interpreted as representative of every month or every COVID-19 phase. The common-month design improves comparability but does not eliminate regime dependence.

\paragraph{Single-source publication selection.}
CafeF is one news source, and its editorial selection is not random. The presence, timing, and wording of an article may depend on firm size, liquidity, market attention, or an ongoing price movement. The news-day control cannot fully remove this selection process.

\paragraph{Entity-linking measurement error.}
The strict threshold improves precision but reduces coverage. The broad candidate file may contain false positives from uppercase abbreviations, while strict filtering may exclude valid ticker mentions. A linking error can connect sentiment to the wrong price series or remove a relevant event entirely. The 32 failed price symbols and 43 strict rows without usable price windows make this limitation observable rather than hypothetical.

\paragraph{Class imbalance and model uncertainty.}
Across the 839 strict predictions, 759 are neutral, 54 are positive, and 26 are negative. Small positive and negative groups produce wide intervals and make group comparisons sensitive to individual observations. Moreover, the current outputs are hard labels without calibrated probabilities, so uncertainty in the classifier is not propagated into the financial tests.

\paragraph{Daily timing and return measurement.}
Daily prices do not reveal whether an article was published before or after the market had already incorporated the information. Mapping every article to the next trading session is conservative and reproducible, but it may delay some events or miss same-day reactions. The price-source adjustment and corporate-action history are also not independently audited in this dataset.

\paragraph{Sparse panel and covariance estimation.}
The period-specific panels have different numbers of ticker--day observations and dates. Sparse fixed-effects designs can produce high leverage and unstable standard errors. The HC1 fallback after a singular HC3 correction is documented, but it reduces the strength of any regression-based inference. The coefficients should therefore be treated as exploratory conditional associations.

\paragraph{Non-causal design and multiple testing.}
The analysis is observational and examines many specifications. It does not control for every confounder, does not use an exogenous news shock, and does not pre-register all hypotheses. Statistical significance in one period or window is not sufficient to establish a causal effect or a deployable trading signal.

\subsection{Future Work}
\label{subsec:future-work-statistical}

Future work should proceed in a pre-specified sequence. First, create an independently annotated and adjudicated CafeF gold set with multiple annotators, report inter-annotator agreement, and evaluate accuracy, macro-F1, balanced accuracy, MCC, and class-specific confusion patterns. The gold set should be split by article or time period so that repeated articles do not leak across train and test data.

Second, repeat inference for all four periods using a frozen model checkpoint and preserve per-class probabilities or logits. This will enable confidence-weighted aggregation, Expected Calibration Error, reliability diagrams, and sensitivity analysis that separates high-confidence from low-confidence predictions.

Third, expand the time axis beyond one month per year, ideally using multiple months before, during, and after the pandemic regime. A longer panel would increase power, allow year and sector fixed effects, and support formal interactions between sentiment and market regime. The broad-linking run should be completed independently and compared with the strict run using pre-defined matching rules.

Fourth, strengthen the financial design by estimating alternative market-adjustment models, including market-adjusted returns and factor models, while fixing the estimation-window rule before inspecting results. Placebo dates, non-news dates, sector decomposition, and clustered standard errors should be included as confirmatory robustness checks rather than selected after seeing the main results.

Finally, if return prediction is a research objective, train a separate chronological forecasting model with a censor gap between the information date and evaluation window. Performance should be compared with simple baselines using out-of-time RMSE, MAE, rank correlation, and economically interpretable evaluation metrics. Any trading simulation would require transaction costs, liquidity constraints, delisting treatment, and strict protection against look-ahead bias.

\subsection{Overall Assessment and Conclusion}
\label{subsec:conclusion-statistical}

The completed analysis establishes a reproducible end-to-end connection between Vietnamese financial headlines and market outcomes. It validates the prediction contract for 839 strict article--ticker inputs across October 2019--2022, constructs target-specific sentiment signals, aligns publication dates to the next trading session, estimates security-level market models, and applies complementary event-study, control-group, group-comparison, panel, and lagged procedures.

The financial evidence is mixed and period-dependent. Aggregate same-day CAR is not statistically different from zero in any year. Negative three-session CAR appears in 2020 and 2021 at nominal significance levels, while the 2021 positive--neutral contrast and the 2020 panel sentiment coefficient are also nominally significant. However, the signs and magnitudes are not stable across years, event windows, or lags; the significant findings arise within a large family of tests; and the confidence intervals and small sentiment groups indicate substantial uncertainty. The news-day control does not reach the conventional 5\% threshold in any period.

The most defensible conclusion is therefore that the pipeline is technically complete and produces observable sentiment--return associations, but the current evidence does not demonstrate a stable, causal, or reliably predictive effect of FinABSA sentiment on Vietnamese stock returns. The occasional significant estimates are useful hypotheses for future confirmatory research, not evidence of a standalone investment strategy.

\begin{table}[htbp]
\centering
\small
\caption{Reproducibility checklist for the completed multi-year validation.}
\label{tab:reproducibility-multi-year}
\begin{tabular}{p{0.42\textwidth}p{0.46\textwidth}}
\toprule
\textbf{Component} & \textbf{Completed implementation} \\
\midrule
Sampling design & October 2019, October 2020, October 2021, October 2022 \\
Price window & May--December of the same year for every period \\
Prediction alignment & Stable \texttt{sample\_id}; 839/839 matched \\
Label normalization & Canonical positive/neutral/negative; zero unknown labels \\
Entity-linking primary sample & Strict confidence threshold $\geq 0.95$ \\
Event mapping & First trading session strictly after publication date \\
Market adjustment & Security-specific market model with VNINDEX proxy \\
Event windows & $[0,0]$, $[0,+1]$, and $[0,+3]$ \\
Inference tests & One-sample $t$-test, sign-flip permutation, bootstrap CI \\
Control test & News-day versus eligible no-news-day comparison \\
Group tests & Welch, permutation, and bootstrap contrasts \\
Panel model & Ticker/date fixed effects; HC1 fallback documented \\
Lag sensitivity & 0, 1, 2, 3, 5, and 10 trading-session lags \\
Output artifacts & CSV tables, period-specific event files, six multi-year figures, Markdown report, Excel workbook \\
\bottomrule
\end{tabular}
\end{table}

The numerical source files are stored under \texttt{data/cafef\_multiyear/analysis\_multiyear/}. The principal tables are \texttt{car\_primary\_all.csv}, \texttt{news\_day\_control\_tests.csv}, \texttt{sentiment\_group\_tests.csv}, \texttt{panel\_key\_coefficients.csv}, \texttt{robustness\_summary.csv}, and \texttt{prediction\_validation\_clean.csv}. The complete workbook is \texttt{multiyear\_validation\_tables.xlsx}, and the period-specific diagnostics are retained under each period's \texttt{analysis/} directory.

\begin{quote}
This is research and analysis only, not personalized financial advice.
\end{quote}
.
%
% Required packages in the main preamble:
% amsmath, amssymb, booktabs, graphicx, longtable, array, hyperref
%
% Figures are expected under images/.
% All numerical values below are from the completed October 2019--2022 runs.
% ============================================================================

\section{Empirical Validation and Statistical Testing}
\label{sec:empirical-validation}

This section evaluates whether target-specific sentiment predictions generated by the FinABSA pipeline are associated with subsequent abnormal stock returns in Vietnamese financial news. The analysis extends the original October 2022 design to four comparable October samples: October 2019, October 2020, October 2021, and October 2022. October 2019 is retained as a pre-pandemic reference period, while 2020--2022 capture different stages of the COVID-19 and post-COVID market environment. The same calendar month is used in every year to reduce month-of-year confounding. Each period is analysed separately using its own price history from May through December of the same year.

The empirical objective is deliberately narrower than a causal claim. The tests ask whether predicted sentiment is statistically associated with market-adjusted price movements after entity linking, time alignment, and aggregation to the ticker--day level. They do not establish that a news article caused a return, that the model generates a profitable trading strategy, or that the observed association generalises to all Vietnamese financial news.

\subsection{Research Questions and Hypotheses}
\label{subsec:research-questions-statistics}

The downstream validation addresses three questions. First, are the model outputs complete and correctly aligned with the target-masked article--ticker inputs? Second, do abnormal returns differ from zero around the first trading session after publication? Third, do the results remain stable across sentiment groups, control samples, fixed-effects specifications, and lag definitions?

The primary null hypotheses are:
\begin{itemize}
    \item $H_{0,1}$: the mean cumulative abnormal return is zero within each event window;
    \item $H_{0,2}$: abnormal returns do not differ between eligible ticker--days with an identified news event and eligible ticker--days without one;
    \item $H_{0,3}$: abnormal returns do not differ systematically across positive, neutral, and negative sentiment groups;
    \item $H_{0,4}$: the conditional association between sentiment and abnormal return is zero after ticker and date fixed effects;
    \item $H_{0,5}$: any observed association does not persist in a stable and monotonic pattern across alternative lags.
\end{itemize}

All hypothesis tests are two-sided and use a nominal significance level of $\alpha=0.05$. Because many period--window--group combinations are examined, individual $p$-values are treated as exploratory unless the pattern is replicated across specifications and survives an appropriate multiple-comparison correction.

\subsection{Data Construction and Multi-year Coverage}
\label{subsec:multi-year-data}

\paragraph{News corpus.}
The crawler collected historical CafeF articles from the October sitemap partitions for each year. The resulting corpus contains 5,133 article records. A stable article identifier is derived from the URL, and each record retains the title, publication timestamp, URL, parsing status, and related metadata. The crawl status is kept explicitly so that an unavailable or malformed page is not silently interpreted as an empty financial article.

A headline can mention more than one company. Since FinABSA predicts sentiment for a specified target rather than discovering ticker symbols itself, the inference unit is an article--ticker pair. Entity linking produces a broad candidate file and a strict high-confidence file. The broad file retains all automatically detected candidates, while the strict file retains candidates with entity-linking confidence at least 0.95. This distinction is a measurement-quality choice rather than a model change.

\begin{table}[htbp]
\centering
\small
\caption{Multi-year CafeF corpus and model-input coverage.}
\label{tab:multi-year-news-coverage}
\begin{tabular}{lrrrrrr}
\toprule
\textbf{Period} & \textbf{Articles} & \textbf{Crawl OK} & \textbf{Full pairs} & \textbf{Strict pairs} & \textbf{Strict + price} & \textbf{Price errors} \\
\midrule
October 2019 & 1,040 & 1,039 & 583 & 154 & 141 & 11 \\
October 2020 & 1,115 & 1,115 & 657 & 182 & 168 & 10 \\
October 2021 & 1,506 & 1,491 & 991 & 297 & 281 & 11 \\
October 2022 & 1,472 & 1,457 & 509 & 206 & 206 & 0 \\
\midrule
\textbf{Total} & \textbf{5,133} & -- & \textbf{2,740} & \textbf{839} & \textbf{796} & \textbf{32} \\
\bottomrule
\end{tabular}
\end{table}

The price errors are not replaced by synthetic observations. They correspond to candidate strings that could not be retrieved as valid securities or to failed endpoint requests. Similarly, the 43 strict input rows without a usable price window are preserved in the inference file but excluded from financial calculations that require a matched price series. This separation prevents missing market observations from being treated as zero returns.

\paragraph{Target masking and prediction contract.}
For each article--ticker pair, the target surface form is replaced with \texttt{Target} and other detected entities are replaced with \texttt{Other}. The model therefore receives a target-masked headline and produces one of three canonical labels: positive, neutral, or negative. The model does not perform named entity recognition or ticker prediction; entity linking is a separate upstream step.

\begin{table}[htbp]
\centering
\small
\caption{Prediction validation by October period.}
\label{tab:multi-year-prediction-validation}
\begin{tabular}{lrrrrrrrr}
\toprule
\textbf{Period} & \textbf{Input} & \textbf{Pred.} & \textbf{Missing} & \textbf{Unexpected} & \textbf{Dup.} & \textbf{Unknown} & \textbf{Neutral} & \textbf{Pos./Neg.} \\
\midrule
October 2019 & 154 & 154 & 0 & 0 & 0 & 0 & 144 & 7 / 3 \\
October 2020 & 182 & 182 & 0 & 0 & 0 & 0 & 167 & 13 / 2 \\
October 2021 & 297 & 297 & 0 & 0 & 0 & 0 & 270 & 16 / 11 \\
October 2022 & 206 & 206 & 0 & 0 & 0 & 0 & 178 & 18 / 10 \\
\midrule
\textbf{Total} & \textbf{839} & \textbf{839} & \textbf{0} & \textbf{0} & \textbf{0} & \textbf{0} & \textbf{759} & \textbf{54 / 26} \\
\bottomrule
\end{tabular}
\end{table}

The prediction file is aligned to the input contract by the stable \texttt{sample\_id}, not by row order. All 839 strict inputs have exactly one matching prediction. There are no missing predictions, unexpected prediction IDs, duplicate IDs, or unknown labels after canonicalisation. This validates the data plumbing, but it does not validate semantic accuracy: without independently annotated gold labels, the prediction distribution cannot be converted into accuracy, macro-F1, MCC, or Cohen's kappa.

\begin{figure}[htbp]
\centering
\includegraphics[width=0.78\textwidth]{images/multiyear_label_distribution.png}
\caption{Composition of the predicted sentiment labels by October period. Neutral predictions dominate every period, while positive and negative groups are comparatively small. The imbalance reduces the precision of group-level comparisons and motivates reporting group sizes together with every estimate.}
\label{fig:multi-year-label-distribution}
\end{figure}

\paragraph{Market data.}
Daily OHLCV data were retrieved through \texttt{vnstock} for the identified securities and the VNINDEX and VN30 reference series. For every October period, the stored price window begins in May of the same year and ends in December. This provides pre-event observations for the market model and post-event observations for lag calculations without importing prices from a different calendar regime.

\begin{table}[htbp]
\centering
\small
\caption{Market-price coverage by October period.}
\label{tab:multi-year-price-coverage}
\begin{tabular}{lrrrrr}
\toprule
\textbf{Period} & \textbf{Price rows} & \textbf{Symbols} & \textbf{Price start} & \textbf{Price end} & \textbf{Strict rows with price} \\
\midrule
October 2019 & 8,320 & 49 & 2019-05-02 & 2019-12-31 & 141 \\
October 2020 & 9,449 & 57 & 2020-05-04 & 2020-12-31 & 168 \\
October 2021 & 11,935 & 73 & 2021-05-04 & 2021-12-31 & 281 \\
October 2022 & 9,182 & 55 & 2022-05-04 & 2022-12-30 & 206 \\
\bottomrule
\end{tabular}
\end{table}

The number of symbols differs by period because the news corpus and entity-linking output differ across years. The price data are therefore comparable in construction, but not identical in cross-sectional composition. This point is important when comparing the magnitude of estimates across periods.

\subsection{Time Alignment and Return Construction}
\label{subsec:time-alignment-multi-year}

The data contain daily prices but not the exact time at which an article was disseminated to market participants. The implemented alignment rule maps each article's normalized publication date to the first trading date strictly after that calendar date. Thus, an article published on a trading day is evaluated against the next available trading session; an article published on a weekend or holiday is also mapped to the next trading session. No fixed 15:00 cutoff is imposed, because the daily data do not identify the intraday time at which the information was incorporated into the price.

For security $i$ on trading date $t$, the continuously compounded percentage return is computed as
\begin{equation}
 r_{i,t}=100\times\left[\log(P_{i,t})-\log(P_{i,t-1})\right],
\end{equation}
where $P_{i,t}$ is the recorded closing price. The market return is computed using the same transformation for VNINDEX. The analysis uses the downloaded price series as stored; any interpretation of return magnitude is conditional on the source's adjustment and corporate-action coverage.

The sentiment score is a deterministic label encoding:
\begin{equation}
 S_{i,t}=\begin{cases}
 +1,&\text{positive prediction},\\
 0,&\text{neutral prediction},\\
 -1,&\text{negative prediction}.
 \end{cases}
\end{equation}

Several article--ticker rows can refer to the same ticker on the same date. Because the market supplies one return observation per ticker--day, the financial unit is formed by aggregation before inference. The panel records mean sentiment, positive/neutral/negative shares, the number of unique articles, and the number of contributing article--ticker rows. This avoids artificially multiplying the effective sample size by repeating one return for every article.

\subsection{Event-Study Methodology}
\label{subsec:event-study-multi-year}

Following the standard market-model event-study framework of Brown and Warner~\cite{brownwarner1985}, the expected return for each security is estimated from observations before 1 October of the relevant year:
\begin{equation}
 R_{i,t}=\alpha_i+\beta_iR_{m,t}+\varepsilon_{i,t},
\end{equation}
where $R_{m,t}$ is the VNINDEX return. The abnormal return is
\begin{equation}
 AR_{i,t}=R_{i,t}-\left(\widehat{\alpha}_i+\widehat{\beta}_iR_{m,t}\right).
\end{equation}
A security-specific market model is fitted when at least 20 pre-event observations are available; otherwise the documented fallback parameters are retained. The final period-specific runs contain 49, 57, 73, and 55 stored price symbols, respectively, with 54 fitted market-model records in the October 2022 run and corresponding period-specific records in the other runs.

For an event aligned to trading date $t$, cumulative abnormal return is defined as
\begin{equation}
 CAR_i(T_1,T_2)=\sum_{\tau=T_1}^{T_2}AR_{i,t+\tau}.
\end{equation}
Three windows are examined: $[0,0]$ for the event session, $[0,+1]$ for the event session and the next session, and $[0,+3]$ for the event session and the next three sessions. For every period and window, the pipeline reports the mean and median CAR, a two-sided one-sample $t$-test against zero, a two-sided sign-flip permutation test with 5,000 replications, and a percentile bootstrap 95\% confidence interval with 5,000 replications. Fixed random seeds make the resampling results reproducible.

\begin{figure}[htbp]
\centering
\includegraphics[width=0.84\textwidth]{images/multiyear_car_by_period.png}
\caption{Aggregate CAR by October period and event window. Points show mean CAR and whiskers show percentile bootstrap 95\% confidence intervals. The figure is intended for cross-period comparison; it does not apply a multiple-comparison correction and does not establish causality.}
\label{fig:multi-year-car}
\end{figure}

\begin{longtable}{llrrrrrr}
\caption{Aggregate event-study CAR tests by period and event window.}\label{tab:multi-year-car-tests}\\
\toprule
\textbf{Period} & \textbf{Window} & \textbf{N} & \textbf{Mean} & \textbf{Median} & \textbf{$p_t$} & \textbf{$p_{perm}$} & \textbf{Bootstrap 95\% CI} \\
\midrule
\endfirsthead
\toprule
\textbf{Period} & \textbf{Window} & \textbf{N} & \textbf{Mean} & \textbf{Median} & \textbf{$p_t$} & \textbf{$p_{perm}$} & \textbf{Bootstrap 95\% CI} \\
\midrule
\endhead
\midrule
\multicolumn{8}{r}{Continued on next page}\\
\endfoot
\bottomrule
\endlastfoot
2019-10 & $[0,0]$ & 120 & -0.1323 & -0.1570 & 0.5420 & 0.5491 & $[-0.5550,0.2916]$ \\
2019-10 & $[0,+1]$ & 120 & 0.0573 & -0.1285 & 0.8833 & 0.8830 & $[-0.6841,0.8108]$ \\
2019-10 & $[0,+3]$ & 120 & 0.0020 & -0.2812 & 0.9970 & 0.9974 & $[-1.0491,1.0923]$ \\
2020-10 & $[0,0]$ & 134 & 0.0522 & -0.2528 & 0.8031 & 0.8072 & $[-0.3645,0.4595]$ \\
2020-10 & $[0,+1]$ & 134 & -0.2975 & -0.6585 & 0.3424 & 0.3363 & $[-0.8991,0.3050]$ \\
2020-10 & $[0,+3]$ & 133 & -0.9261 & -1.3113 & 0.0438 & 0.0440 & $[-1.8254,-0.0657]$ \\
2021-10 & $[0,0]$ & 212 & -0.2797 & -0.4321 & 0.1070 & 0.1096 & $[-0.6117,0.0587]$ \\
2021-10 & $[0,+1]$ & 212 & -0.2584 & -0.5174 & 0.2920 & 0.2953 & $[-0.7377,0.2377]$ \\
2021-10 & $[0,+3]$ & 212 & -1.0149 & -1.0412 & 0.0101 & 0.0088 & $[-1.7967,-0.2662]$ \\
2022-10 & $[0,0]$ & 150 & 0.1363 & 0.0008 & 0.6023 & 0.6091 & $[-0.3557,0.6421]$ \\
2022-10 & $[0,+1]$ & 150 & -0.0220 & 0.1798 & 0.9501 & 0.9562 & $[-0.7298,0.6360]$ \\
2022-10 & $[0,+3]$ & 150 & -0.2889 & -0.7977 & 0.5614 & 0.5615 & $[-1.2725,0.6793]$ \\
\end{longtable}

\paragraph{Event-study results.}
The aggregate CAR results do not support a single stable sentiment-linked price response across the four October periods. The same-day window is statistically indistinguishable from zero in every period: all $p$-values exceed 0.10 and all bootstrap intervals include zero. In the three-session window, 2020-10 and 2021-10 produce negative mean CAR estimates with nominal $p$-values below 0.05. These two results are noteworthy but should not be treated as decisive evidence: they are two findings among many period--window tests, they are not adjusted for multiple comparisons, and the negative three-session CAR is not itself evidence that sentiment caused the decline.

The 2019 and 2022 estimates remain close to zero or statistically imprecise across all windows. This cross-period heterogeneity is more consistent with a period-dependent market association than with a universal FinABSA effect. It may reflect differences in market regime, news composition, entity coverage, macroeconomic shocks, or the composition of the available price universe.

\subsection{News-Day Versus No-News-Day Control}
\label{subsec:news-control-multi-year}

The event study is complemented by a control comparison. For each period, the control sample contains eligible ticker--days belonging to tickers that appear in the event sample, restricted to October of that year. A news ticker--day is a ticker--day with at least one aligned article event; a no-news ticker--day is an eligible ticker--day without an identified event. The estimand is
\begin{equation}
 \Delta=\overline{AR}_{news}-\overline{AR}_{no\text{-}news}.
\end{equation}

The two distributions are compared using a Welch unequal-variance test, a pooled-label permutation test with 5,000 permutations, and a percentile bootstrap interval for the difference in means.

\begin{table}[htbp]
\centering
\scriptsize
\caption{News-day versus eligible no-news-day abnormal returns.}
\label{tab:news-control-multi-year}
\begin{tabular}{lrrrrrrrr}
\toprule
\textbf{Period} & \textbf{$N_{news}$} & \textbf{$N_{0}$} & \textbf{News mean} & \textbf{No-news mean} & \textbf{$\Delta$} & \textbf{Welch $p$} & \textbf{Perm. $p$} & \textbf{Bootstrap 95\% CI} \\
\midrule
2019-10 & 120 & 941 & -0.1323 & 0.0064 & -0.1386 & 0.5383 & 0.4585 & $[-0.5937,0.3095]$ \\
2020-10 & 134 & 1,027 & 0.0522 & -0.2161 & 0.2683 & 0.2240 & 0.1828 & $[-0.1601,0.7062]$ \\
2021-10 & 212 & 1,217 & -0.2797 & -0.1039 & -0.1758 & 0.3553 & 0.3771 & $[-0.5493,0.1980]$ \\
2022-10 & 150 & 936 & 0.1363 & -0.3559 & 0.4921 & 0.0804 & 0.0684 & $[-0.0645,1.0442]$ \\
\bottomrule
\end{tabular}
\end{table}

\begin{figure}[htbp]
\centering
\includegraphics[width=0.78\textwidth]{images/multiyear_news_day_control.png}
\caption{News-day minus eligible no-news-day abnormal return by October period. Whiskers show percentile bootstrap 95\% confidence intervals; each interval includes zero.}
\label{fig:multi-year-news-control}
\end{figure}

The control comparison does not reject equality at the 5\% level in any period. The largest positive point estimate occurs in October 2022, where news days have 0.4921 percentage points higher abnormal returns than no-news days, but the Welch and permutation $p$-values are 0.0804 and 0.0684, respectively, and the bootstrap interval crosses zero. The result is therefore suggestive but not conventionally significant. The comparison remains observational because news selection is not random: firms that receive coverage may systematically differ from firms that do not.

\subsection{Sentiment-Group Contrasts}
\label{subsec:sentiment-groups-multi-year}

At the ticker--day level, the analysis compares three pre-specified contrasts: positive minus neutral, negative minus neutral, and positive minus negative. Each contrast is evaluated with a Welch test, a pooled permutation test, and a bootstrap interval. Differences are defined as group A minus group B.

\begin{longtable}{llrrrrrr}
\caption{Sentiment-group abnormal-return contrasts by period.}\label{tab:sentiment-contrasts-multi-year}\\
\toprule
\textbf{Period} & \textbf{Contrast} & \textbf{$N_A$} & \textbf{$N_B$} & \textbf{Difference} & \textbf{Welch $p$} & \textbf{Perm. $p$} & \textbf{Bootstrap 95\% CI} \\
\midrule
\endfirsthead
\toprule
\textbf{Period} & \textbf{Contrast} & \textbf{$N_A$} & \textbf{$N_B$} & \textbf{Difference} & \textbf{Welch $p$} & \textbf{Perm. $p$} & \textbf{Bootstrap 95\% CI} \\
\midrule
\endhead
\midrule
\multicolumn{8}{r}{Continued on next page}\\
\endfoot
\bottomrule
\endlastfoot
2019-10 & Positive $-$ Neutral & 6 & 113 & -1.5118 & 0.1196 & 0.1252 & $[-3.2365,-0.3238]$ \\
2019-10 & Negative $-$ Neutral & 1 & 113 & -- & -- & -- & Not estimable \\
2019-10 & Positive $-$ Negative & 6 & 1 & -- & -- & -- & Not estimable \\
2020-10 & Positive $-$ Neutral & 10 & 123 & -0.8544 & 0.3380 & 0.2813 & $[-2.1566,0.9352]$ \\
2020-10 & Negative $-$ Neutral & 1 & 123 & -- & -- & -- & Not estimable \\
2020-10 & Positive $-$ Negative & 10 & 1 & -- & -- & -- & Not estimable \\
2021-10 & Positive $-$ Neutral & 12 & 191 & 1.5693 & 0.0267 & 0.0388 & $[0.5210,2.8516]$ \\
2021-10 & Negative $-$ Neutral & 9 & 191 & 0.0915 & 0.9093 & 0.9112 & $[-1.0946,1.7607]$ \\
2021-10 & Positive $-$ Negative & 12 & 9 & 1.4779 & 0.1460 & 0.1372 & $[-0.4072,3.2104]$ \\
2022-10 & Positive $-$ Neutral & 14 & 126 & 0.6990 & 0.4518 & 0.4271 & $[-0.9308,2.4970]$ \\
2022-10 & Negative $-$ Neutral & 10 & 126 & -0.4688 & 0.6712 & 0.6601 & $[-2.5243,1.4876]$ \\
2022-10 & Positive $-$ Negative & 14 & 10 & 1.1678 & 0.3962 & 0.3971 & $[-1.2938,3.6937]$ \\
\end{longtable}

The most notable group contrast appears in October 2021: the positive group has a 1.5693 percentage-point higher mean abnormal return than the neutral group, with Welch $p=0.0267$, permutation $p=0.0388$, and a bootstrap interval that excludes zero. This result should be described as a period-specific exploratory finding rather than a general law. The positive group contains only 12 ticker--day observations, the negative group contains nine, the analysis involves multiple comparisons, and the effect is not replicated in the corresponding positive--neutral comparisons for 2019, 2020, or 2022.

The 2019 positive--neutral point estimate is negative, whereas the 2020 and 2022 differences are also statistically imprecise. The failure of sign consistency across years weakens the claim that predicted sentiment produces a stable cross-sectional ranking of abnormal returns.

\begin{figure}[htbp]
\centering
\includegraphics[width=0.78\textwidth]{images/multiyear_positive_neutral_contrast.png}
\caption{Positive-minus-neutral abnormal-return contrast by October period. Whiskers show bootstrap 95\% confidence intervals. The positive 2021 estimate is not sufficient on its own to establish a stable sentiment effect across years.}
\label{fig:multi-year-positive-neutral}
\end{figure}

\subsection{Two-Way Fixed-Effects Panel Regression}
\label{subsec:panel-regression-multi-year}

To control for persistent differences between securities and common shocks across dates, the following two-way fixed-effects model is estimated separately for each October period:
\begin{equation}
 AR_{i,\tau(i,t)}=\alpha_i+\delta_{\tau(i,t)}+\beta_1S_{i,t}
 +\beta_2\mathrm{NegativeShare}_{i,t}
 +\beta_3\mathrm{PositiveShare}_{i,t}
 +\gamma\log(1+\mathrm{ArticleCount}_{i,t})+\varepsilon_{i,t},
\end{equation}
where $\alpha_i$ denotes ticker fixed effects, $\delta_{\tau(i,t)}$ denotes date fixed effects, and $\tau(i,t)$ is the first trading session strictly after publication. The dependent variable is market-model abnormal return on the aligned evaluation date.

\begin{table}[htbp]
\centering
\scriptsize
\caption{Key coefficients from period-specific two-way fixed-effects regressions. The full coefficient tables, including ticker and date fixed-effect indicators, are provided in the accompanying CSV output.}
\label{tab:panel-key-coefficients-multi-year}
\begin{tabular}{llrrrr}
\toprule
\textbf{Period} & \textbf{Variable} & \textbf{Coefficient} & \textbf{Std. error} & \textbf{$t$-statistic} & \textbf{$p$-value} \\
\midrule
2019-10 & Sentiment score & -0.4450 & 0.8508 & -0.5230 & 0.6009 \\
2019-10 & Negative share & 1.5084 & 1.8412 & 0.8192 & 0.4127 \\
2019-10 & Positive share & 1.0633 & 1.2562 & 0.8468 & 0.3973 \\
2019-10 & Log article count & -1.5405 & 3.3930 & -0.4540 & 0.6498 \\
2020-10 & Sentiment score & 2.5056 & 0.8042 & 3.1156 & 0.0018 \\
2020-10 & Negative share & -4.9332 & 1.8033 & -2.7356 & 0.0062 \\
2020-10 & Positive share & -2.4276 & 1.2770 & -1.9010 & 0.0573 \\
2020-10 & Log article count & 2.3897 & 1.9986 & 1.1957 & 0.2318 \\
2021-10 & Sentiment score & -0.9696 & 2.4681 & -0.3927 & 0.6944 \\
2021-10 & Negative share & 0.1799 & 2.5174 & 0.0715 & 0.9430 \\
2021-10 & Positive share & 0.4345 & 2.6128 & 0.1663 & 0.8679 \\
2021-10 & Log article count & 0.5867 & 1.0124 & 0.5795 & 0.5623 \\
2022-10 & Sentiment score & 6.5062 & 5.3763 & 1.2100 & 0.2262 \\
2022-10 & Negative share & 5.9827 & 5.2919 & 1.1305 & 0.2582 \\
2022-10 & Positive share & -6.1476 & 5.5086 & -1.1160 & 0.2644 \\
2022-10 & Log article count & 0.3079 & 1.3849 & 0.2224 & 0.8241 \\
\bottomrule
\end{tabular}
\end{table}

The sentiment-score coefficient is positive and statistically significant only in the 2020 period. The coefficient is negative and insignificant in 2019 and 2021, and positive but imprecise in 2022. This sign instability is central: a robust sentiment--return association would be expected to exhibit a more consistent direction under comparable definitions. The 2020 coefficient therefore constitutes a period-specific association, not evidence that the model has a stable return-prediction relationship.

The coefficient on negative share is also significant only in 2020 and is negative, while the positive-share coefficient is not significant at the 5\% level. Because sentiment shares are compositional and mechanically related to the omitted neutral share, their coefficients should not be interpreted as independent causal marginal effects. The relatively high within-period $R^2$ values are not evidence of strong predictive power: ticker and date fixed effects can absorb substantial variation, while the sentiment covariates remain sensitive to sample composition.

The intended covariance estimator was HC3. In sparse two-way fixed-effects designs, HC3 can produce infinite standard errors when leverage reaches a singular boundary. The implementation therefore falls back to HC1 when necessary and records this fallback in the period-specific metadata. These panel coefficients should consequently be treated as exploratory conditional associations.

\begin{figure}[htbp]
\centering
\includegraphics[width=0.78\textwidth]{images/multiyear_panel_sentiment_coefficient.png}
\caption{Period-specific sentiment-score coefficients from the two-way fixed-effects models. Whiskers show coefficient plus or minus 1.96 standard errors. The 2020 coefficient is positive and relatively precise, but the sign is not replicated across the other periods.}
\label{fig:multi-year-panel-coefficient}
\end{figure}

\subsection{Lagged Robustness and Sensitivity Analysis}
\label{subsec:lag-robustness-multi-year}

Lag analysis evaluates abnormal returns at 0, 1, 2, 3, 5, and 10 trading sessions after the aligned event date. These are sensitivity checks rather than independent causal tests. The sample size decreases at longer lags because the relevant security price must remain available after the event.

\begin{longtable}{llrrrrrr}
\caption{Lagged abnormal-return robustness by period.}\label{tab:lag-robustness-multi-year}\\
\toprule
\textbf{Period} & \textbf{Lag} & \textbf{N} & \textbf{Overall AR} & \textbf{Negative} & \textbf{Neutral} & \textbf{Positive} \\
\midrule
\endfirsthead
\toprule
\textbf{Period} & \textbf{Lag} & \textbf{N} & \textbf{Overall AR} & \textbf{Negative} & \textbf{Neutral} & \textbf{Positive} \\
\midrule
\endhead
\midrule
\multicolumn{7}{r}{Continued on next page}\\
\endfoot
\bottomrule
\endlastfoot
2019-10 & 0 & 120 & -0.1323 & -1.5766 & -0.0648 & 0.8940 \\
2019-10 & 1 & 146 & 0.3194 & 0.0588 & 0.2772 & 1.2785 \\
2019-10 & 2 & 141 & -0.0968 & 0.0331 & -0.0974 & -0.1286 \\
2019-10 & 3 & 134 & 0.0067 & -0.4507 & 0.0283 & -0.3248 \\
2019-10 & 5 & 108 & 0.1154 & -0.6470 & 0.1166 & 0.2232 \\
2019-10 & 10 & 74 & -0.1635 & 0.1766 & -0.1547 & -0.6273 \\
2020-10 & 0 & 134 & 0.2261 & -1.9519 & 0.3085 & -0.4498 \\
2020-10 & 1 & 164 & -0.1520 & 0.5470 & -0.1975 & 0.2915 \\
2020-10 & 2 & 159 & -0.3772 & 1.3165 & -0.3780 & -0.7051 \\
2020-10 & 3 & 152 & 0.0687 & 0.9315 & 0.1112 & -0.6345 \\
2020-10 & 5 & 144 & -0.0997 & 0.4069 & -0.1472 & 0.4869 \\
2020-10 & 10 & 100 & -0.4924 & 0.3574 & -0.5596 & 1.5471 \\
2021-10 & 0 & 212 & -0.0285 & -0.1815 & -0.0939 & 1.0708 \\
2021-10 & 1 & 267 & 0.2065 & 0.6397 & 0.2368 & -0.5095 \\
2021-10 & 2 & 255 & -0.3156 & -0.4267 & -0.2764 & -0.8200 \\
2021-10 & 3 & 244 & 0.0773 & -0.3897 & 0.1940 & -1.2377 \\
2021-10 & 5 & 219 & 0.1398 & -0.1027 & 0.1966 & -0.4387 \\
2021-10 & 10 & 147 & -0.0471 & 1.2318 & -0.1101 & -0.1871 \\
2022-10 & 0 & 150 & -0.0936 & -0.7884 & -0.1653 & 0.8940 \\
2022-10 & 1 & 174 & -0.1797 & 0.3368 & 0.0713 & -2.7331 \\
2022-10 & 2 & 163 & -0.2007 & -1.2344 & -0.0991 & -0.4067 \\
2022-10 & 3 & 153 & -0.5972 & -2.6554 & -0.4837 & -0.1760 \\
2022-10 & 5 & 134 & -0.7243 & 0.4479 & -0.9416 & 0.5251 \\
2022-10 & 10 & 93 & -0.5183 & -0.4422 & -0.2779 & -4.1530 \\
\end{longtable}

\begin{figure}[htbp]
\centering
\includegraphics[width=0.84\textwidth]{images/multiyear_lag_robustness.png}
\caption{Lagged sensitivity of mean abnormal returns across the four October periods. The changing signs and magnitudes indicate that the association is sensitive to the evaluation horizon rather than a stable delayed response.}
\label{fig:multi-year-lag-robustness}
\end{figure}

The lag results do not display a common monotonic decay. In 2019, the overall estimate changes from negative at lag 0 to positive at lag 1 and back toward zero thereafter. In 2020, the estimates become more negative at lags 2 and 10, while in 2021 the sign alternates across horizons. In 2022, the overall mean is increasingly negative at lags 3 and 5. Sentiment-specific estimates are even less stable, particularly for the small positive and negative groups. The appropriate conclusion is that the apparent price association depends on the timing convention and the available evaluation horizon.

\subsection{Additional Diagnostic Evidence}
\label{subsec:additional-diagnostics-multi-year}

The output package also contains daily ticker--day sentiment and abnormal-return tables, market-model coverage, event observations, and panel-regression data. These artifacts support auditability and allow the reader to trace each reported result back to an aligned ticker--day observation. The following diagnostic figures are descriptive rather than inferential.

\begin{figure}[htbp]
\centering
\includegraphics[width=0.82\textwidth]{images/market_model_coverage.png}
\caption{Market-model coverage diagnostic from the completed October 2022 run. It shows the number of pre-event observations available for individual symbols and illustrates why coverage thresholds and fallback handling are necessary.}
\label{fig:market-model-coverage-multi-year}
\end{figure}

\begin{figure}[htbp]
\centering
\includegraphics[width=0.90\textwidth]{images/sentiment_vs_abnormal_return.png}
\caption{Descriptive comparison between daily mean sentiment and mean abnormal return in the ticker--day panel. The figure is not a causal identification plot and does not control for all market-wide shocks.}
\label{fig:sentiment-vs-ar-multi-year}
\end{figure}

\subsection{Multiple Comparisons and Statistical Interpretation}
\label{subsec:multiple-comparisons}

The multi-year design produces several families of related tests: four periods, three event windows, multiple sentiment groups, a news-day control, panel coefficients, and six lag horizons. Consequently, at least some nominally small $p$-values can occur by chance when many hypotheses are examined. The results are therefore interpreted using three criteria: the size and direction of the estimate, whether the confidence interval excludes zero, and whether the finding is replicated under related specifications or periods.

Under this framework, the 2020 and 2021 three-day CAR findings and the 2021 positive--neutral contrast are interesting signals for further study, but not definitive discoveries. They should be reported with their exact sample sizes, intervals, and unadjusted $p$-values. A confirmatory study should pre-specify a smaller number of primary hypotheses and apply a family-wise or false-discovery-rate correction, such as Holm or Benjamini--Hochberg, depending on the inferential objective.

\subsection{Analyses Not Validly Estimated in the Current Run}
\label{subsec:deferred-analyses-multi-year}

Several analyses appeared in the initial experimental plan but are not assigned numerical claims here because the required artifacts are absent. First, intrinsic CafeF accuracy, macro-F1, MCC, and Cohen's kappa require an independently human-annotated gold-label sample. The model's own predictions cannot serve as their ground truth. Second, probability calibration metrics such as Expected Calibration Error require valid per-class probabilities or logits; the current prediction files contain hard labels and empty probability fields. Third, out-of-time return forecasting requires a pre-specified chronological training, validation, and test design with a fitted return-prediction model; the current pipeline is an event-study association analysis, not a trained forecaster. Fourth, a broad-linking robustness result requires a complete independent inference and financial analysis run on the 2,740-row broad input set; strict and full input files alone are not sufficient.

This separation is methodologically important. It prevents unsupported accuracy, calibration, forecasting, and robustness numbers from being presented as results. The completed outputs can be extended later without changing the core data contract.

\subsection{Limitations}
\label{subsec:limitations-statistical}

\paragraph{Limited temporal coverage and regime dependence.}
Although the study now covers four October periods, it still samples only one month per year. October 2019--2022 may contain unusual news intensity, market stress, or sector composition. The estimates should not be interpreted as representative of every month or every COVID-19 phase. The common-month design improves comparability but does not eliminate regime dependence.

\paragraph{Single-source publication selection.}
CafeF is one news source, and its editorial selection is not random. The presence, timing, and wording of an article may depend on firm size, liquidity, market attention, or an ongoing price movement. The news-day control cannot fully remove this selection process.

\paragraph{Entity-linking measurement error.}
The strict threshold improves precision but reduces coverage. The broad candidate file may contain false positives from uppercase abbreviations, while strict filtering may exclude valid ticker mentions. A linking error can connect sentiment to the wrong price series or remove a relevant event entirely. The 32 failed price symbols and 43 strict rows without usable price windows make this limitation observable rather than hypothetical.

\paragraph{Class imbalance and model uncertainty.}
Across the 839 strict predictions, 759 are neutral, 54 are positive, and 26 are negative. Small positive and negative groups produce wide intervals and make group comparisons sensitive to individual observations. Moreover, the current outputs are hard labels without calibrated probabilities, so uncertainty in the classifier is not propagated into the financial tests.

\paragraph{Daily timing and return measurement.}
Daily prices do not reveal whether an article was published before or after the market had already incorporated the information. Mapping every article to the next trading session is conservative and reproducible, but it may delay some events or miss same-day reactions. The price-source adjustment and corporate-action history are also not independently audited in this dataset.

\paragraph{Sparse panel and covariance estimation.}
The period-specific panels have different numbers of ticker--day observations and dates. Sparse fixed-effects designs can produce high leverage and unstable standard errors. The HC1 fallback after a singular HC3 correction is documented, but it reduces the strength of any regression-based inference. The coefficients should therefore be treated as exploratory conditional associations.

\paragraph{Non-causal design and multiple testing.}
The analysis is observational and examines many specifications. It does not control for every confounder, does not use an exogenous news shock, and does not pre-register all hypotheses. Statistical significance in one period or window is not sufficient to establish a causal effect or a deployable trading signal.

\subsection{Future Work}
\label{subsec:future-work-statistical}

Future work should proceed in a pre-specified sequence. First, create an independently annotated and adjudicated CafeF gold set with multiple annotators, report inter-annotator agreement, and evaluate accuracy, macro-F1, balanced accuracy, MCC, and class-specific confusion patterns. The gold set should be split by article or time period so that repeated articles do not leak across train and test data.

Second, repeat inference for all four periods using a frozen model checkpoint and preserve per-class probabilities or logits. This will enable confidence-weighted aggregation, Expected Calibration Error, reliability diagrams, and sensitivity analysis that separates high-confidence from low-confidence predictions.

Third, expand the time axis beyond one month per year, ideally using multiple months before, during, and after the pandemic regime. A longer panel would increase power, allow year and sector fixed effects, and support formal interactions between sentiment and market regime. The broad-linking run should be completed independently and compared with the strict run using pre-defined matching rules.

Fourth, strengthen the financial design by estimating alternative market-adjustment models, including market-adjusted returns and factor models, while fixing the estimation-window rule before inspecting results. Placebo dates, non-news dates, sector decomposition, and clustered standard errors should be included as confirmatory robustness checks rather than selected after seeing the main results.

Finally, if return prediction is a research objective, train a separate chronological forecasting model with a censor gap between the information date and evaluation window. Performance should be compared with simple baselines using out-of-time RMSE, MAE, rank correlation, and economically interpretable evaluation metrics. Any trading simulation would require transaction costs, liquidity constraints, delisting treatment, and strict protection against look-ahead bias.

\subsection{Overall Assessment and Conclusion}
\label{subsec:conclusion-statistical}

The completed analysis establishes a reproducible end-to-end connection between Vietnamese financial headlines and market outcomes. It validates the prediction contract for 839 strict article--ticker inputs across October 2019--2022, constructs target-specific sentiment signals, aligns publication dates to the next trading session, estimates security-level market models, and applies complementary event-study, control-group, group-comparison, panel, and lagged procedures.

The financial evidence is mixed and period-dependent. Aggregate same-day CAR is not statistically different from zero in any year. Negative three-session CAR appears in 2020 and 2021 at nominal significance levels, while the 2021 positive--neutral contrast and the 2020 panel sentiment coefficient are also nominally significant. However, the signs and magnitudes are not stable across years, event windows, or lags; the significant findings arise within a large family of tests; and the confidence intervals and small sentiment groups indicate substantial uncertainty. The news-day control does not reach the conventional 5\% threshold in any period.

The most defensible conclusion is therefore that the pipeline is technically complete and produces observable sentiment--return associations, but the current evidence does not demonstrate a stable, causal, or reliably predictive effect of FinABSA sentiment on Vietnamese stock returns. The occasional significant estimates are useful hypotheses for future confirmatory research, not evidence of a standalone investment strategy.

\begin{table}[htbp]
\centering
\small
\caption{Reproducibility checklist for the completed multi-year validation.}
\label{tab:reproducibility-multi-year}
\begin{tabular}{p{0.42\textwidth}p{0.46\textwidth}}
\toprule
\textbf{Component} & \textbf{Completed implementation} \\
\midrule
Sampling design & October 2019, October 2020, October 2021, October 2022 \\
Price window & May--December of the same year for every period \\
Prediction alignment & Stable \texttt{sample\_id}; 839/839 matched \\
Label normalization & Canonical positive/neutral/negative; zero unknown labels \\
Entity-linking primary sample & Strict confidence threshold $\geq 0.95$ \\
Event mapping & First trading session strictly after publication date \\
Market adjustment & Security-specific market model with VNINDEX proxy \\
Event windows & $[0,0]$, $[0,+1]$, and $[0,+3]$ \\
Inference tests & One-sample $t$-test, sign-flip permutation, bootstrap CI \\
Control test & News-day versus eligible no-news-day comparison \\
Group tests & Welch, permutation, and bootstrap contrasts \\
Panel model & Ticker/date fixed effects; HC1 fallback documented \\
Lag sensitivity & 0, 1, 2, 3, 5, and 10 trading-session lags \\
Output artifacts & CSV tables, period-specific event files, six multi-year figures, Markdown report, Excel workbook \\
\bottomrule
\end{tabular}
\end{table}

The numerical source files are stored under \texttt{data/cafef\_multiyear/analysis\_multiyear/}. The principal tables are \texttt{car\_primary\_all.csv}, \texttt{news\_day\_control\_tests.csv}, \texttt{sentiment\_group\_tests.csv}, \texttt{panel\_key\_coefficients.csv}, \texttt{robustness\_summary.csv}, and \texttt{prediction\_validation\_clean.csv}. The complete workbook is \texttt{multiyear\_validation\_tables.xlsx}, and the period-specific diagnostics are retained under each period's \texttt{analysis/} directory.

\begin{quote}
This is research and analysis only, not personalized financial advice.
\end{quote}
.
%
% Required packages in the main preamble:
% amsmath, amssymb, booktabs, graphicx, longtable, array, hyperref
%
% Figures are expected under images/.
% All numerical values below are from the completed October 2019--2022 runs.
% ============================================================================

\section{Empirical Validation and Statistical Testing}
\label{sec:empirical-validation}

This section evaluates whether target-specific sentiment predictions generated by the FinABSA pipeline are associated with subsequent abnormal stock returns in Vietnamese financial news. The analysis extends the original October 2022 design to four comparable October samples: October 2019, October 2020, October 2021, and October 2022. October 2019 is retained as a pre-pandemic reference period, while 2020--2022 capture different stages of the COVID-19 and post-COVID market environment. The same calendar month is used in every year to reduce month-of-year confounding. Each period is analysed separately using its own price history from May through December of the same year.

The empirical objective is deliberately narrower than a causal claim. The tests ask whether predicted sentiment is statistically associated with market-adjusted price movements after entity linking, time alignment, and aggregation to the ticker--day level. They do not establish that a news article caused a return, that the model generates a profitable trading strategy, or that the observed association generalises to all Vietnamese financial news.

\subsection{Research Questions and Hypotheses}
\label{subsec:research-questions-statistics}

The downstream validation addresses three questions. First, are the model outputs complete and correctly aligned with the target-masked article--ticker inputs? Second, do abnormal returns differ from zero around the first trading session after publication? Third, do the results remain stable across sentiment groups, control samples, fixed-effects specifications, and lag definitions?

The primary null hypotheses are:
\begin{itemize}
    \item $H_{0,1}$: the mean cumulative abnormal return is zero within each event window;
    \item $H_{0,2}$: abnormal returns do not differ between eligible ticker--days with an identified news event and eligible ticker--days without one;
    \item $H_{0,3}$: abnormal returns do not differ systematically across positive, neutral, and negative sentiment groups;
    \item $H_{0,4}$: the conditional association between sentiment and abnormal return is zero after ticker and date fixed effects;
    \item $H_{0,5}$: any observed association does not persist in a stable and monotonic pattern across alternative lags.
\end{itemize}

All hypothesis tests are two-sided and use a nominal significance level of $\alpha=0.05$. Because many period--window--group combinations are examined, individual $p$-values are treated as exploratory unless the pattern is replicated across specifications and survives an appropriate multiple-comparison correction.

\subsection{Data Construction and Multi-year Coverage}
\label{subsec:multi-year-data}

\paragraph{News corpus.}
The crawler collected historical CafeF articles from the October sitemap partitions for each year. The resulting corpus contains 5,133 article records. A stable article identifier is derived from the URL, and each record retains the title, publication timestamp, URL, parsing status, and related metadata. The crawl status is kept explicitly so that an unavailable or malformed page is not silently interpreted as an empty financial article.

A headline can mention more than one company. Since FinABSA predicts sentiment for a specified target rather than discovering ticker symbols itself, the inference unit is an article--ticker pair. Entity linking produces a broad candidate file and a strict high-confidence file. The broad file retains all automatically detected candidates, while the strict file retains candidates with entity-linking confidence at least 0.95. This distinction is a measurement-quality choice rather than a model change.

\begin{table}[htbp]
\centering
\small
\caption{Multi-year CafeF corpus and model-input coverage.}
\label{tab:multi-year-news-coverage}
\begin{tabular}{lrrrrrr}
\toprule
\textbf{Period} & \textbf{Articles} & \textbf{Crawl OK} & \textbf{Full pairs} & \textbf{Strict pairs} & \textbf{Strict + price} & \textbf{Price errors} \\
\midrule
October 2019 & 1,040 & 1,039 & 583 & 154 & 141 & 11 \\
October 2020 & 1,115 & 1,115 & 657 & 182 & 168 & 10 \\
October 2021 & 1,506 & 1,491 & 991 & 297 & 281 & 11 \\
October 2022 & 1,472 & 1,457 & 509 & 206 & 206 & 0 \\
\midrule
\textbf{Total} & \textbf{5,133} & -- & \textbf{2,740} & \textbf{839} & \textbf{796} & \textbf{32} \\
\bottomrule
\end{tabular}
\end{table}

The price errors are not replaced by synthetic observations. They correspond to candidate strings that could not be retrieved as valid securities or to failed endpoint requests. Similarly, the 43 strict input rows without a usable price window are preserved in the inference file but excluded from financial calculations that require a matched price series. This separation prevents missing market observations from being treated as zero returns.

\paragraph{Target masking and prediction contract.}
For each article--ticker pair, the target surface form is replaced with \texttt{Target} and other detected entities are replaced with \texttt{Other}. The model therefore receives a target-masked headline and produces one of three canonical labels: positive, neutral, or negative. The model does not perform named entity recognition or ticker prediction; entity linking is a separate upstream step.

\begin{table}[htbp]
\centering
\small
\caption{Prediction validation by October period.}
\label{tab:multi-year-prediction-validation}
\begin{tabular}{lrrrrrrrr}
\toprule
\textbf{Period} & \textbf{Input} & \textbf{Pred.} & \textbf{Missing} & \textbf{Unexpected} & \textbf{Dup.} & \textbf{Unknown} & \textbf{Neutral} & \textbf{Pos./Neg.} \\
\midrule
October 2019 & 154 & 154 & 0 & 0 & 0 & 0 & 144 & 7 / 3 \\
October 2020 & 182 & 182 & 0 & 0 & 0 & 0 & 167 & 13 / 2 \\
October 2021 & 297 & 297 & 0 & 0 & 0 & 0 & 270 & 16 / 11 \\
October 2022 & 206 & 206 & 0 & 0 & 0 & 0 & 178 & 18 / 10 \\
\midrule
\textbf{Total} & \textbf{839} & \textbf{839} & \textbf{0} & \textbf{0} & \textbf{0} & \textbf{0} & \textbf{759} & \textbf{54 / 26} \\
\bottomrule
\end{tabular}
\end{table}

The prediction file is aligned to the input contract by the stable \texttt{sample\_id}, not by row order. All 839 strict inputs have exactly one matching prediction. There are no missing predictions, unexpected prediction IDs, duplicate IDs, or unknown labels after canonicalisation. This validates the data plumbing, but it does not validate semantic accuracy: without independently annotated gold labels, the prediction distribution cannot be converted into accuracy, macro-F1, MCC, or Cohen's kappa.

\begin{figure}[htbp]
\centering
\includegraphics[width=0.78\textwidth]{images/multiyear_label_distribution.png}
\caption{Composition of the predicted sentiment labels by October period. Neutral predictions dominate every period, while positive and negative groups are comparatively small. The imbalance reduces the precision of group-level comparisons and motivates reporting group sizes together with every estimate.}
\label{fig:multi-year-label-distribution}
\end{figure}

\paragraph{Market data.}
Daily OHLCV data were retrieved through \texttt{vnstock} for the identified securities and the VNINDEX and VN30 reference series. For every October period, the stored price window begins in May of the same year and ends in December. This provides pre-event observations for the market model and post-event observations for lag calculations without importing prices from a different calendar regime.

\begin{table}[htbp]
\centering
\small
\caption{Market-price coverage by October period.}
\label{tab:multi-year-price-coverage}
\begin{tabular}{lrrrrr}
\toprule
\textbf{Period} & \textbf{Price rows} & \textbf{Symbols} & \textbf{Price start} & \textbf{Price end} & \textbf{Strict rows with price} \\
\midrule
October 2019 & 8,320 & 49 & 2019-05-02 & 2019-12-31 & 141 \\
October 2020 & 9,449 & 57 & 2020-05-04 & 2020-12-31 & 168 \\
October 2021 & 11,935 & 73 & 2021-05-04 & 2021-12-31 & 281 \\
October 2022 & 9,182 & 55 & 2022-05-04 & 2022-12-30 & 206 \\
\bottomrule
\end{tabular}
\end{table}

The number of symbols differs by period because the news corpus and entity-linking output differ across years. The price data are therefore comparable in construction, but not identical in cross-sectional composition. This point is important when comparing the magnitude of estimates across periods.

\subsection{Time Alignment and Return Construction}
\label{subsec:time-alignment-multi-year}

The data contain daily prices but not the exact time at which an article was disseminated to market participants. The implemented alignment rule maps each article's normalized publication date to the first trading date strictly after that calendar date. Thus, an article published on a trading day is evaluated against the next available trading session; an article published on a weekend or holiday is also mapped to the next trading session. No fixed 15:00 cutoff is imposed, because the daily data do not identify the intraday time at which the information was incorporated into the price.

For security $i$ on trading date $t$, the continuously compounded percentage return is computed as
\begin{equation}
 r_{i,t}=100\times\left[\log(P_{i,t})-\log(P_{i,t-1})\right],
\end{equation}
where $P_{i,t}$ is the recorded closing price. The market return is computed using the same transformation for VNINDEX. The analysis uses the downloaded price series as stored; any interpretation of return magnitude is conditional on the source's adjustment and corporate-action coverage.

The sentiment score is a deterministic label encoding:
\begin{equation}
 S_{i,t}=\begin{cases}
 +1,&\text{positive prediction},\\
 0,&\text{neutral prediction},\\
 -1,&\text{negative prediction}.
 \end{cases}
\end{equation}

Several article--ticker rows can refer to the same ticker on the same date. Because the market supplies one return observation per ticker--day, the financial unit is formed by aggregation before inference. The panel records mean sentiment, positive/neutral/negative shares, the number of unique articles, and the number of contributing article--ticker rows. This avoids artificially multiplying the effective sample size by repeating one return for every article.

\subsection{Event-Study Methodology}
\label{subsec:event-study-multi-year}

Following the standard market-model event-study framework of Brown and Warner~\cite{brownwarner1985}, the expected return for each security is estimated from observations before 1 October of the relevant year:
\begin{equation}
 R_{i,t}=\alpha_i+\beta_iR_{m,t}+\varepsilon_{i,t},
\end{equation}
where $R_{m,t}$ is the VNINDEX return. The abnormal return is
\begin{equation}
 AR_{i,t}=R_{i,t}-\left(\widehat{\alpha}_i+\widehat{\beta}_iR_{m,t}\right).
\end{equation}
A security-specific market model is fitted when at least 20 pre-event observations are available; otherwise the documented fallback parameters are retained. The final period-specific runs contain 49, 57, 73, and 55 stored price symbols, respectively, with 54 fitted market-model records in the October 2022 run and corresponding period-specific records in the other runs.

For an event aligned to trading date $t$, cumulative abnormal return is defined as
\begin{equation}
 CAR_i(T_1,T_2)=\sum_{\tau=T_1}^{T_2}AR_{i,t+\tau}.
\end{equation}
Three windows are examined: $[0,0]$ for the event session, $[0,+1]$ for the event session and the next session, and $[0,+3]$ for the event session and the next three sessions. For every period and window, the pipeline reports the mean and median CAR, a two-sided one-sample $t$-test against zero, a two-sided sign-flip permutation test with 5,000 replications, and a percentile bootstrap 95\% confidence interval with 5,000 replications. Fixed random seeds make the resampling results reproducible.

\begin{figure}[htbp]
\centering
\includegraphics[width=0.84\textwidth]{images/multiyear_car_by_period.png}
\caption{Aggregate CAR by October period and event window. Points show mean CAR and whiskers show percentile bootstrap 95\% confidence intervals. The figure is intended for cross-period comparison; it does not apply a multiple-comparison correction and does not establish causality.}
\label{fig:multi-year-car}
\end{figure}

\begin{longtable}{llrrrrrr}
\caption{Aggregate event-study CAR tests by period and event window.}\label{tab:multi-year-car-tests}\\
\toprule
\textbf{Period} & \textbf{Window} & \textbf{N} & \textbf{Mean} & \textbf{Median} & \textbf{$p_t$} & \textbf{$p_{perm}$} & \textbf{Bootstrap 95\% CI} \\
\midrule
\endfirsthead
\toprule
\textbf{Period} & \textbf{Window} & \textbf{N} & \textbf{Mean} & \textbf{Median} & \textbf{$p_t$} & \textbf{$p_{perm}$} & \textbf{Bootstrap 95\% CI} \\
\midrule
\endhead
\midrule
\multicolumn{8}{r}{Continued on next page}\\
\endfoot
\bottomrule
\endlastfoot
2019-10 & $[0,0]$ & 120 & -0.1323 & -0.1570 & 0.5420 & 0.5491 & $[-0.5550,0.2916]$ \\
2019-10 & $[0,+1]$ & 120 & 0.0573 & -0.1285 & 0.8833 & 0.8830 & $[-0.6841,0.8108]$ \\
2019-10 & $[0,+3]$ & 120 & 0.0020 & -0.2812 & 0.9970 & 0.9974 & $[-1.0491,1.0923]$ \\
2020-10 & $[0,0]$ & 134 & 0.0522 & -0.2528 & 0.8031 & 0.8072 & $[-0.3645,0.4595]$ \\
2020-10 & $[0,+1]$ & 134 & -0.2975 & -0.6585 & 0.3424 & 0.3363 & $[-0.8991,0.3050]$ \\
2020-10 & $[0,+3]$ & 133 & -0.9261 & -1.3113 & 0.0438 & 0.0440 & $[-1.8254,-0.0657]$ \\
2021-10 & $[0,0]$ & 212 & -0.2797 & -0.4321 & 0.1070 & 0.1096 & $[-0.6117,0.0587]$ \\
2021-10 & $[0,+1]$ & 212 & -0.2584 & -0.5174 & 0.2920 & 0.2953 & $[-0.7377,0.2377]$ \\
2021-10 & $[0,+3]$ & 212 & -1.0149 & -1.0412 & 0.0101 & 0.0088 & $[-1.7967,-0.2662]$ \\
2022-10 & $[0,0]$ & 150 & 0.1363 & 0.0008 & 0.6023 & 0.6091 & $[-0.3557,0.6421]$ \\
2022-10 & $[0,+1]$ & 150 & -0.0220 & 0.1798 & 0.9501 & 0.9562 & $[-0.7298,0.6360]$ \\
2022-10 & $[0,+3]$ & 150 & -0.2889 & -0.7977 & 0.5614 & 0.5615 & $[-1.2725,0.6793]$ \\
\end{longtable}

\paragraph{Event-study results.}
The aggregate CAR results do not support a single stable sentiment-linked price response across the four October periods. The same-day window is statistically indistinguishable from zero in every period: all $p$-values exceed 0.10 and all bootstrap intervals include zero. In the three-session window, 2020-10 and 2021-10 produce negative mean CAR estimates with nominal $p$-values below 0.05. These two results are noteworthy but should not be treated as decisive evidence: they are two findings among many period--window tests, they are not adjusted for multiple comparisons, and the negative three-session CAR is not itself evidence that sentiment caused the decline.

The 2019 and 2022 estimates remain close to zero or statistically imprecise across all windows. This cross-period heterogeneity is more consistent with a period-dependent market association than with a universal FinABSA effect. It may reflect differences in market regime, news composition, entity coverage, macroeconomic shocks, or the composition of the available price universe.

\subsection{News-Day Versus No-News-Day Control}
\label{subsec:news-control-multi-year}

The event study is complemented by a control comparison. For each period, the control sample contains eligible ticker--days belonging to tickers that appear in the event sample, restricted to October of that year. A news ticker--day is a ticker--day with at least one aligned article event; a no-news ticker--day is an eligible ticker--day without an identified event. The estimand is
\begin{equation}
 \Delta=\overline{AR}_{news}-\overline{AR}_{no\text{-}news}.
\end{equation}

The two distributions are compared using a Welch unequal-variance test, a pooled-label permutation test with 5,000 permutations, and a percentile bootstrap interval for the difference in means.

\begin{table}[htbp]
\centering
\scriptsize
\caption{News-day versus eligible no-news-day abnormal returns.}
\label{tab:news-control-multi-year}
\begin{tabular}{lrrrrrrrr}
\toprule
\textbf{Period} & \textbf{$N_{news}$} & \textbf{$N_{0}$} & \textbf{News mean} & \textbf{No-news mean} & \textbf{$\Delta$} & \textbf{Welch $p$} & \textbf{Perm. $p$} & \textbf{Bootstrap 95\% CI} \\
\midrule
2019-10 & 120 & 941 & -0.1323 & 0.0064 & -0.1386 & 0.5383 & 0.4585 & $[-0.5937,0.3095]$ \\
2020-10 & 134 & 1,027 & 0.0522 & -0.2161 & 0.2683 & 0.2240 & 0.1828 & $[-0.1601,0.7062]$ \\
2021-10 & 212 & 1,217 & -0.2797 & -0.1039 & -0.1758 & 0.3553 & 0.3771 & $[-0.5493,0.1980]$ \\
2022-10 & 150 & 936 & 0.1363 & -0.3559 & 0.4921 & 0.0804 & 0.0684 & $[-0.0645,1.0442]$ \\
\bottomrule
\end{tabular}
\end{table}

\begin{figure}[htbp]
\centering
\includegraphics[width=0.78\textwidth]{images/multiyear_news_day_control.png}
\caption{News-day minus eligible no-news-day abnormal return by October period. Whiskers show percentile bootstrap 95\% confidence intervals; each interval includes zero.}
\label{fig:multi-year-news-control}
\end{figure}

The control comparison does not reject equality at the 5\% level in any period. The largest positive point estimate occurs in October 2022, where news days have 0.4921 percentage points higher abnormal returns than no-news days, but the Welch and permutation $p$-values are 0.0804 and 0.0684, respectively, and the bootstrap interval crosses zero. The result is therefore suggestive but not conventionally significant. The comparison remains observational because news selection is not random: firms that receive coverage may systematically differ from firms that do not.

\subsection{Sentiment-Group Contrasts}
\label{subsec:sentiment-groups-multi-year}

At the ticker--day level, the analysis compares three pre-specified contrasts: positive minus neutral, negative minus neutral, and positive minus negative. Each contrast is evaluated with a Welch test, a pooled permutation test, and a bootstrap interval. Differences are defined as group A minus group B.

\begin{longtable}{llrrrrrr}
\caption{Sentiment-group abnormal-return contrasts by period.}\label{tab:sentiment-contrasts-multi-year}\\
\toprule
\textbf{Period} & \textbf{Contrast} & \textbf{$N_A$} & \textbf{$N_B$} & \textbf{Difference} & \textbf{Welch $p$} & \textbf{Perm. $p$} & \textbf{Bootstrap 95\% CI} \\
\midrule
\endfirsthead
\toprule
\textbf{Period} & \textbf{Contrast} & \textbf{$N_A$} & \textbf{$N_B$} & \textbf{Difference} & \textbf{Welch $p$} & \textbf{Perm. $p$} & \textbf{Bootstrap 95\% CI} \\
\midrule
\endhead
\midrule
\multicolumn{8}{r}{Continued on next page}\\
\endfoot
\bottomrule
\endlastfoot
2019-10 & Positive $-$ Neutral & 6 & 113 & -1.5118 & 0.1196 & 0.1252 & $[-3.2365,-0.3238]$ \\
2019-10 & Negative $-$ Neutral & 1 & 113 & -- & -- & -- & Not estimable \\
2019-10 & Positive $-$ Negative & 6 & 1 & -- & -- & -- & Not estimable \\
2020-10 & Positive $-$ Neutral & 10 & 123 & -0.8544 & 0.3380 & 0.2813 & $[-2.1566,0.9352]$ \\
2020-10 & Negative $-$ Neutral & 1 & 123 & -- & -- & -- & Not estimable \\
2020-10 & Positive $-$ Negative & 10 & 1 & -- & -- & -- & Not estimable \\
2021-10 & Positive $-$ Neutral & 12 & 191 & 1.5693 & 0.0267 & 0.0388 & $[0.5210,2.8516]$ \\
2021-10 & Negative $-$ Neutral & 9 & 191 & 0.0915 & 0.9093 & 0.9112 & $[-1.0946,1.7607]$ \\
2021-10 & Positive $-$ Negative & 12 & 9 & 1.4779 & 0.1460 & 0.1372 & $[-0.4072,3.2104]$ \\
2022-10 & Positive $-$ Neutral & 14 & 126 & 0.6990 & 0.4518 & 0.4271 & $[-0.9308,2.4970]$ \\
2022-10 & Negative $-$ Neutral & 10 & 126 & -0.4688 & 0.6712 & 0.6601 & $[-2.5243,1.4876]$ \\
2022-10 & Positive $-$ Negative & 14 & 10 & 1.1678 & 0.3962 & 0.3971 & $[-1.2938,3.6937]$ \\
\end{longtable}

The most notable group contrast appears in October 2021: the positive group has a 1.5693 percentage-point higher mean abnormal return than the neutral group, with Welch $p=0.0267$, permutation $p=0.0388$, and a bootstrap interval that excludes zero. This result should be described as a period-specific exploratory finding rather than a general law. The positive group contains only 12 ticker--day observations, the negative group contains nine, the analysis involves multiple comparisons, and the effect is not replicated in the corresponding positive--neutral comparisons for 2019, 2020, or 2022.

The 2019 positive--neutral point estimate is negative, whereas the 2020 and 2022 differences are also statistically imprecise. The failure of sign consistency across years weakens the claim that predicted sentiment produces a stable cross-sectional ranking of abnormal returns.

\begin{figure}[htbp]
\centering
\includegraphics[width=0.78\textwidth]{images/multiyear_positive_neutral_contrast.png}
\caption{Positive-minus-neutral abnormal-return contrast by October period. Whiskers show bootstrap 95\% confidence intervals. The positive 2021 estimate is not sufficient on its own to establish a stable sentiment effect across years.}
\label{fig:multi-year-positive-neutral}
\end{figure}

\subsection{Two-Way Fixed-Effects Panel Regression}
\label{subsec:panel-regression-multi-year}

To control for persistent differences between securities and common shocks across dates, the following two-way fixed-effects model is estimated separately for each October period:
\begin{equation}
 AR_{i,\tau(i,t)}=\alpha_i+\delta_{\tau(i,t)}+\beta_1S_{i,t}
 +\beta_2\mathrm{NegativeShare}_{i,t}
 +\beta_3\mathrm{PositiveShare}_{i,t}
 +\gamma\log(1+\mathrm{ArticleCount}_{i,t})+\varepsilon_{i,t},
\end{equation}
where $\alpha_i$ denotes ticker fixed effects, $\delta_{\tau(i,t)}$ denotes date fixed effects, and $\tau(i,t)$ is the first trading session strictly after publication. The dependent variable is market-model abnormal return on the aligned evaluation date.

\begin{table}[htbp]
\centering
\scriptsize
\caption{Key coefficients from period-specific two-way fixed-effects regressions. The full coefficient tables, including ticker and date fixed-effect indicators, are provided in the accompanying CSV output.}
\label{tab:panel-key-coefficients-multi-year}
\begin{tabular}{llrrrr}
\toprule
\textbf{Period} & \textbf{Variable} & \textbf{Coefficient} & \textbf{Std. error} & \textbf{$t$-statistic} & \textbf{$p$-value} \\
\midrule
2019-10 & Sentiment score & -0.4450 & 0.8508 & -0.5230 & 0.6009 \\
2019-10 & Negative share & 1.5084 & 1.8412 & 0.8192 & 0.4127 \\
2019-10 & Positive share & 1.0633 & 1.2562 & 0.8468 & 0.3973 \\
2019-10 & Log article count & -1.5405 & 3.3930 & -0.4540 & 0.6498 \\
2020-10 & Sentiment score & 2.5056 & 0.8042 & 3.1156 & 0.0018 \\
2020-10 & Negative share & -4.9332 & 1.8033 & -2.7356 & 0.0062 \\
2020-10 & Positive share & -2.4276 & 1.2770 & -1.9010 & 0.0573 \\
2020-10 & Log article count & 2.3897 & 1.9986 & 1.1957 & 0.2318 \\
2021-10 & Sentiment score & -0.9696 & 2.4681 & -0.3927 & 0.6944 \\
2021-10 & Negative share & 0.1799 & 2.5174 & 0.0715 & 0.9430 \\
2021-10 & Positive share & 0.4345 & 2.6128 & 0.1663 & 0.8679 \\
2021-10 & Log article count & 0.5867 & 1.0124 & 0.5795 & 0.5623 \\
2022-10 & Sentiment score & 6.5062 & 5.3763 & 1.2100 & 0.2262 \\
2022-10 & Negative share & 5.9827 & 5.2919 & 1.1305 & 0.2582 \\
2022-10 & Positive share & -6.1476 & 5.5086 & -1.1160 & 0.2644 \\
2022-10 & Log article count & 0.3079 & 1.3849 & 0.2224 & 0.8241 \\
\bottomrule
\end{tabular}
\end{table}

The sentiment-score coefficient is positive and statistically significant only in the 2020 period. The coefficient is negative and insignificant in 2019 and 2021, and positive but imprecise in 2022. This sign instability is central: a robust sentiment--return association would be expected to exhibit a more consistent direction under comparable definitions. The 2020 coefficient therefore constitutes a period-specific association, not evidence that the model has a stable return-prediction relationship.

The coefficient on negative share is also significant only in 2020 and is negative, while the positive-share coefficient is not significant at the 5\% level. Because sentiment shares are compositional and mechanically related to the omitted neutral share, their coefficients should not be interpreted as independent causal marginal effects. The relatively high within-period $R^2$ values are not evidence of strong predictive power: ticker and date fixed effects can absorb substantial variation, while the sentiment covariates remain sensitive to sample composition.

The intended covariance estimator was HC3. In sparse two-way fixed-effects designs, HC3 can produce infinite standard errors when leverage reaches a singular boundary. The implementation therefore falls back to HC1 when necessary and records this fallback in the period-specific metadata. These panel coefficients should consequently be treated as exploratory conditional associations.

\begin{figure}[htbp]
\centering
\includegraphics[width=0.78\textwidth]{images/multiyear_panel_sentiment_coefficient.png}
\caption{Period-specific sentiment-score coefficients from the two-way fixed-effects models. Whiskers show coefficient plus or minus 1.96 standard errors. The 2020 coefficient is positive and relatively precise, but the sign is not replicated across the other periods.}
\label{fig:multi-year-panel-coefficient}
\end{figure}

\subsection{Lagged Robustness and Sensitivity Analysis}
\label{subsec:lag-robustness-multi-year}

Lag analysis evaluates abnormal returns at 0, 1, 2, 3, 5, and 10 trading sessions after the aligned event date. These are sensitivity checks rather than independent causal tests. The sample size decreases at longer lags because the relevant security price must remain available after the event.

\begin{longtable}{llrrrrrr}
\caption{Lagged abnormal-return robustness by period.}\label{tab:lag-robustness-multi-year}\\
\toprule
\textbf{Period} & \textbf{Lag} & \textbf{N} & \textbf{Overall AR} & \textbf{Negative} & \textbf{Neutral} & \textbf{Positive} \\
\midrule
\endfirsthead
\toprule
\textbf{Period} & \textbf{Lag} & \textbf{N} & \textbf{Overall AR} & \textbf{Negative} & \textbf{Neutral} & \textbf{Positive} \\
\midrule
\endhead
\midrule
\multicolumn{7}{r}{Continued on next page}\\
\endfoot
\bottomrule
\endlastfoot
2019-10 & 0 & 120 & -0.1323 & -1.5766 & -0.0648 & 0.8940 \\
2019-10 & 1 & 146 & 0.3194 & 0.0588 & 0.2772 & 1.2785 \\
2019-10 & 2 & 141 & -0.0968 & 0.0331 & -0.0974 & -0.1286 \\
2019-10 & 3 & 134 & 0.0067 & -0.4507 & 0.0283 & -0.3248 \\
2019-10 & 5 & 108 & 0.1154 & -0.6470 & 0.1166 & 0.2232 \\
2019-10 & 10 & 74 & -0.1635 & 0.1766 & -0.1547 & -0.6273 \\
2020-10 & 0 & 134 & 0.2261 & -1.9519 & 0.3085 & -0.4498 \\
2020-10 & 1 & 164 & -0.1520 & 0.5470 & -0.1975 & 0.2915 \\
2020-10 & 2 & 159 & -0.3772 & 1.3165 & -0.3780 & -0.7051 \\
2020-10 & 3 & 152 & 0.0687 & 0.9315 & 0.1112 & -0.6345 \\
2020-10 & 5 & 144 & -0.0997 & 0.4069 & -0.1472 & 0.4869 \\
2020-10 & 10 & 100 & -0.4924 & 0.3574 & -0.5596 & 1.5471 \\
2021-10 & 0 & 212 & -0.0285 & -0.1815 & -0.0939 & 1.0708 \\
2021-10 & 1 & 267 & 0.2065 & 0.6397 & 0.2368 & -0.5095 \\
2021-10 & 2 & 255 & -0.3156 & -0.4267 & -0.2764 & -0.8200 \\
2021-10 & 3 & 244 & 0.0773 & -0.3897 & 0.1940 & -1.2377 \\
2021-10 & 5 & 219 & 0.1398 & -0.1027 & 0.1966 & -0.4387 \\
2021-10 & 10 & 147 & -0.0471 & 1.2318 & -0.1101 & -0.1871 \\
2022-10 & 0 & 150 & -0.0936 & -0.7884 & -0.1653 & 0.8940 \\
2022-10 & 1 & 174 & -0.1797 & 0.3368 & 0.0713 & -2.7331 \\
2022-10 & 2 & 163 & -0.2007 & -1.2344 & -0.0991 & -0.4067 \\
2022-10 & 3 & 153 & -0.5972 & -2.6554 & -0.4837 & -0.1760 \\
2022-10 & 5 & 134 & -0.7243 & 0.4479 & -0.9416 & 0.5251 \\
2022-10 & 10 & 93 & -0.5183 & -0.4422 & -0.2779 & -4.1530 \\
\end{longtable}

\begin{figure}[htbp]
\centering
\includegraphics[width=0.84\textwidth]{images/multiyear_lag_robustness.png}
\caption{Lagged sensitivity of mean abnormal returns across the four October periods. The changing signs and magnitudes indicate that the association is sensitive to the evaluation horizon rather than a stable delayed response.}
\label{fig:multi-year-lag-robustness}
\end{figure}

The lag results do not display a common monotonic decay. In 2019, the overall estimate changes from negative at lag 0 to positive at lag 1 and back toward zero thereafter. In 2020, the estimates become more negative at lags 2 and 10, while in 2021 the sign alternates across horizons. In 2022, the overall mean is increasingly negative at lags 3 and 5. Sentiment-specific estimates are even less stable, particularly for the small positive and negative groups. The appropriate conclusion is that the apparent price association depends on the timing convention and the available evaluation horizon.

\subsection{Additional Diagnostic Evidence}
\label{subsec:additional-diagnostics-multi-year}

The output package also contains daily ticker--day sentiment and abnormal-return tables, market-model coverage, event observations, and panel-regression data. These artifacts support auditability and allow the reader to trace each reported result back to an aligned ticker--day observation. The following diagnostic figures are descriptive rather than inferential.

\begin{figure}[htbp]
\centering
\includegraphics[width=0.82\textwidth]{images/market_model_coverage.png}
\caption{Market-model coverage diagnostic from the completed October 2022 run. It shows the number of pre-event observations available for individual symbols and illustrates why coverage thresholds and fallback handling are necessary.}
\label{fig:market-model-coverage-multi-year}
\end{figure}

\begin{figure}[htbp]
\centering
\includegraphics[width=0.90\textwidth]{images/sentiment_vs_abnormal_return.png}
\caption{Descriptive comparison between daily mean sentiment and mean abnormal return in the ticker--day panel. The figure is not a causal identification plot and does not control for all market-wide shocks.}
\label{fig:sentiment-vs-ar-multi-year}
\end{figure}

\subsection{Multiple Comparisons and Statistical Interpretation}
\label{subsec:multiple-comparisons}

The multi-year design produces several families of related tests: four periods, three event windows, multiple sentiment groups, a news-day control, panel coefficients, and six lag horizons. Consequently, at least some nominally small $p$-values can occur by chance when many hypotheses are examined. The results are therefore interpreted using three criteria: the size and direction of the estimate, whether the confidence interval excludes zero, and whether the finding is replicated under related specifications or periods.

Under this framework, the 2020 and 2021 three-day CAR findings and the 2021 positive--neutral contrast are interesting signals for further study, but not definitive discoveries. They should be reported with their exact sample sizes, intervals, and unadjusted $p$-values. A confirmatory study should pre-specify a smaller number of primary hypotheses and apply a family-wise or false-discovery-rate correction, such as Holm or Benjamini--Hochberg, depending on the inferential objective.

\subsection{Analyses Not Validly Estimated in the Current Run}
\label{subsec:deferred-analyses-multi-year}

Several analyses appeared in the initial experimental plan but are not assigned numerical claims here because the required artifacts are absent. First, intrinsic CafeF accuracy, macro-F1, MCC, and Cohen's kappa require an independently human-annotated gold-label sample. The model's own predictions cannot serve as their ground truth. Second, probability calibration metrics such as Expected Calibration Error require valid per-class probabilities or logits; the current prediction files contain hard labels and empty probability fields. Third, out-of-time return forecasting requires a pre-specified chronological training, validation, and test design with a fitted return-prediction model; the current pipeline is an event-study association analysis, not a trained forecaster. Fourth, a broad-linking robustness result requires a complete independent inference and financial analysis run on the 2,740-row broad input set; strict and full input files alone are not sufficient.

This separation is methodologically important. It prevents unsupported accuracy, calibration, forecasting, and robustness numbers from being presented as results. The completed outputs can be extended later without changing the core data contract.

\subsection{Limitations}
\label{subsec:limitations-statistical}

\paragraph{Limited temporal coverage and regime dependence.}
Although the study now covers four October periods, it still samples only one month per year. October 2019--2022 may contain unusual news intensity, market stress, or sector composition. The estimates should not be interpreted as representative of every month or every COVID-19 phase. The common-month design improves comparability but does not eliminate regime dependence.

\paragraph{Single-source publication selection.}
CafeF is one news source, and its editorial selection is not random. The presence, timing, and wording of an article may depend on firm size, liquidity, market attention, or an ongoing price movement. The news-day control cannot fully remove this selection process.

\paragraph{Entity-linking measurement error.}
The strict threshold improves precision but reduces coverage. The broad candidate file may contain false positives from uppercase abbreviations, while strict filtering may exclude valid ticker mentions. A linking error can connect sentiment to the wrong price series or remove a relevant event entirely. The 32 failed price symbols and 43 strict rows without usable price windows make this limitation observable rather than hypothetical.

\paragraph{Class imbalance and model uncertainty.}
Across the 839 strict predictions, 759 are neutral, 54 are positive, and 26 are negative. Small positive and negative groups produce wide intervals and make group comparisons sensitive to individual observations. Moreover, the current outputs are hard labels without calibrated probabilities, so uncertainty in the classifier is not propagated into the financial tests.

\paragraph{Daily timing and return measurement.}
Daily prices do not reveal whether an article was published before or after the market had already incorporated the information. Mapping every article to the next trading session is conservative and reproducible, but it may delay some events or miss same-day reactions. The price-source adjustment and corporate-action history are also not independently audited in this dataset.

\paragraph{Sparse panel and covariance estimation.}
The period-specific panels have different numbers of ticker--day observations and dates. Sparse fixed-effects designs can produce high leverage and unstable standard errors. The HC1 fallback after a singular HC3 correction is documented, but it reduces the strength of any regression-based inference. The coefficients should therefore be treated as exploratory conditional associations.

\paragraph{Non-causal design and multiple testing.}
The analysis is observational and examines many specifications. It does not control for every confounder, does not use an exogenous news shock, and does not pre-register all hypotheses. Statistical significance in one period or window is not sufficient to establish a causal effect or a deployable trading signal.

\subsection{Future Work}
\label{subsec:future-work-statistical}

Future work should proceed in a pre-specified sequence. First, create an independently annotated and adjudicated CafeF gold set with multiple annotators, report inter-annotator agreement, and evaluate accuracy, macro-F1, balanced accuracy, MCC, and class-specific confusion patterns. The gold set should be split by article or time period so that repeated articles do not leak across train and test data.

Second, repeat inference for all four periods using a frozen model checkpoint and preserve per-class probabilities or logits. This will enable confidence-weighted aggregation, Expected Calibration Error, reliability diagrams, and sensitivity analysis that separates high-confidence from low-confidence predictions.

Third, expand the time axis beyond one month per year, ideally using multiple months before, during, and after the pandemic regime. A longer panel would increase power, allow year and sector fixed effects, and support formal interactions between sentiment and market regime. The broad-linking run should be completed independently and compared with the strict run using pre-defined matching rules.

Fourth, strengthen the financial design by estimating alternative market-adjustment models, including market-adjusted returns and factor models, while fixing the estimation-window rule before inspecting results. Placebo dates, non-news dates, sector decomposition, and clustered standard errors should be included as confirmatory robustness checks rather than selected after seeing the main results.

Finally, if return prediction is a research objective, train a separate chronological forecasting model with a censor gap between the information date and evaluation window. Performance should be compared with simple baselines using out-of-time RMSE, MAE, rank correlation, and economically interpretable evaluation metrics. Any trading simulation would require transaction costs, liquidity constraints, delisting treatment, and strict protection against look-ahead bias.

\subsection{Overall Assessment and Conclusion}
\label{subsec:conclusion-statistical}

The completed analysis establishes a reproducible end-to-end connection between Vietnamese financial headlines and market outcomes. It validates the prediction contract for 839 strict article--ticker inputs across October 2019--2022, constructs target-specific sentiment signals, aligns publication dates to the next trading session, estimates security-level market models, and applies complementary event-study, control-group, group-comparison, panel, and lagged procedures.

The financial evidence is mixed and period-dependent. Aggregate same-day CAR is not statistically different from zero in any year. Negative three-session CAR appears in 2020 and 2021 at nominal significance levels, while the 2021 positive--neutral contrast and the 2020 panel sentiment coefficient are also nominally significant. However, the signs and magnitudes are not stable across years, event windows, or lags; the significant findings arise within a large family of tests; and the confidence intervals and small sentiment groups indicate substantial uncertainty. The news-day control does not reach the conventional 5\% threshold in any period.

The most defensible conclusion is therefore that the pipeline is technically complete and produces observable sentiment--return associations, but the current evidence does not demonstrate a stable, causal, or reliably predictive effect of FinABSA sentiment on Vietnamese stock returns. The occasional significant estimates are useful hypotheses for future confirmatory research, not evidence of a standalone investment strategy.

\begin{table}[htbp]
\centering
\small
\caption{Reproducibility checklist for the completed multi-year validation.}
\label{tab:reproducibility-multi-year}
\begin{tabular}{p{0.42\textwidth}p{0.46\textwidth}}
\toprule
\textbf{Component} & \textbf{Completed implementation} \\
\midrule
Sampling design & October 2019, October 2020, October 2021, October 2022 \\
Price window & May--December of the same year for every period \\
Prediction alignment & Stable \texttt{sample\_id}; 839/839 matched \\
Label normalization & Canonical positive/neutral/negative; zero unknown labels \\
Entity-linking primary sample & Strict confidence threshold $\geq 0.95$ \\
Event mapping & First trading session strictly after publication date \\
Market adjustment & Security-specific market model with VNINDEX proxy \\
Event windows & $[0,0]$, $[0,+1]$, and $[0,+3]$ \\
Inference tests & One-sample $t$-test, sign-flip permutation, bootstrap CI \\
Control test & News-day versus eligible no-news-day comparison \\
Group tests & Welch, permutation, and bootstrap contrasts \\
Panel model & Ticker/date fixed effects; HC1 fallback documented \\
Lag sensitivity & 0, 1, 2, 3, 5, and 10 trading-session lags \\
Output artifacts & CSV tables, period-specific event files, six multi-year figures, Markdown report, Excel workbook \\
\bottomrule
\end{tabular}
\end{table}

The numerical source files are stored under \texttt{data/cafef\_multiyear/analysis\_multiyear/}. The principal tables are \texttt{car\_primary\_all.csv}, \texttt{news\_day\_control\_tests.csv}, \texttt{sentiment\_group\_tests.csv}, \texttt{panel\_key\_coefficients.csv}, \texttt{robustness\_summary.csv}, and \texttt{prediction\_validation\_clean.csv}. The complete workbook is \texttt{multiyear\_validation\_tables.xlsx}, and the period-specific diagnostics are retained under each period's \texttt{analysis/} directory.

\begin{quote}
This is research and analysis only, not personalized financial advice.
\end{quote}
.
%
% Required packages in the main preamble:
% amsmath, amssymb, booktabs, graphicx, longtable, array, hyperref
%
% Figures are expected under images/.
% All numerical values below are from the completed October 2019--2022 runs.
% ============================================================================

\section{Empirical Validation and Statistical Testing}
\label{sec:empirical-validation}

This section evaluates whether target-specific sentiment predictions generated by the FinABSA pipeline are associated with subsequent abnormal stock returns in Vietnamese financial news. The analysis extends the original October 2022 design to four comparable October samples: October 2019, October 2020, October 2021, and October 2022. October 2019 is retained as a pre-pandemic reference period, while 2020--2022 capture different stages of the COVID-19 and post-COVID market environment. The same calendar month is used in every year to reduce month-of-year confounding. Each period is analysed separately using its own price history from May through December of the same year.

The empirical objective is deliberately narrower than a causal claim. The tests ask whether predicted sentiment is statistically associated with market-adjusted price movements after entity linking, time alignment, and aggregation to the ticker--day level. They do not establish that a news article caused a return, that the model generates a profitable trading strategy, or that the observed association generalises to all Vietnamese financial news.

\subsection{Research Questions and Hypotheses}
\label{subsec:research-questions-statistics}

The downstream validation addresses three questions. First, are the model outputs complete and correctly aligned with the target-masked article--ticker inputs? Second, do abnormal returns differ from zero around the first trading session after publication? Third, do the results remain stable across sentiment groups, control samples, fixed-effects specifications, and lag definitions?

The primary null hypotheses are:
\begin{itemize}
    \item $H_{0,1}$: the mean cumulative abnormal return is zero within each event window;
    \item $H_{0,2}$: abnormal returns do not differ between eligible ticker--days with an identified news event and eligible ticker--days without one;
    \item $H_{0,3}$: abnormal returns do not differ systematically across positive, neutral, and negative sentiment groups;
    \item $H_{0,4}$: the conditional association between sentiment and abnormal return is zero after ticker and date fixed effects;
    \item $H_{0,5}$: any observed association does not persist in a stable and monotonic pattern across alternative lags.
\end{itemize}

All hypothesis tests are two-sided and use a nominal significance level of $\alpha=0.05$. Because many period--window--group combinations are examined, individual $p$-values are treated as exploratory unless the pattern is replicated across specifications and survives an appropriate multiple-comparison correction.

\subsection{Data Construction and Multi-year Coverage}
\label{subsec:multi-year-data}

\paragraph{News corpus.}
The crawler collected historical CafeF articles from the October sitemap partitions for each year. The resulting corpus contains 5,133 article records. A stable article identifier is derived from the URL, and each record retains the title, publication timestamp, URL, parsing status, and related metadata. The crawl status is kept explicitly so that an unavailable or malformed page is not silently interpreted as an empty financial article.

A headline can mention more than one company. Since FinABSA predicts sentiment for a specified target rather than discovering ticker symbols itself, the inference unit is an article--ticker pair. Entity linking produces a broad candidate file and a strict high-confidence file. The broad file retains all automatically detected candidates, while the strict file retains candidates with entity-linking confidence at least 0.95. This distinction is a measurement-quality choice rather than a model change.

\begin{table}[htbp]
\centering
\small
\caption{Multi-year CafeF corpus and model-input coverage.}
\label{tab:multi-year-news-coverage}
\begin{tabular}{lrrrrrr}
\toprule
\textbf{Period} & \textbf{Articles} & \textbf{Crawl OK} & \textbf{Full pairs} & \textbf{Strict pairs} & \textbf{Strict + price} & \textbf{Price errors} \\
\midrule
October 2019 & 1,040 & 1,039 & 583 & 154 & 141 & 11 \\
October 2020 & 1,115 & 1,115 & 657 & 182 & 168 & 10 \\
October 2021 & 1,506 & 1,491 & 991 & 297 & 281 & 11 \\
October 2022 & 1,472 & 1,457 & 509 & 206 & 206 & 0 \\
\midrule
\textbf{Total} & \textbf{5,133} & -- & \textbf{2,740} & \textbf{839} & \textbf{796} & \textbf{32} \\
\bottomrule
\end{tabular}
\end{table}

The price errors are not replaced by synthetic observations. They correspond to candidate strings that could not be retrieved as valid securities or to failed endpoint requests. Similarly, the 43 strict input rows without a usable price window are preserved in the inference file but excluded from financial calculations that require a matched price series. This separation prevents missing market observations from being treated as zero returns.

\paragraph{Target masking and prediction contract.}
For each article--ticker pair, the target surface form is replaced with \texttt{Target} and other detected entities are replaced with \texttt{Other}. The model therefore receives a target-masked headline and produces one of three canonical labels: positive, neutral, or negative. The model does not perform named entity recognition or ticker prediction; entity linking is a separate upstream step.

\begin{table}[htbp]
\centering
\small
\caption{Prediction validation by October period.}
\label{tab:multi-year-prediction-validation}
\begin{tabular}{lrrrrrrrr}
\toprule
\textbf{Period} & \textbf{Input} & \textbf{Pred.} & \textbf{Missing} & \textbf{Unexpected} & \textbf{Dup.} & \textbf{Unknown} & \textbf{Neutral} & \textbf{Pos./Neg.} \\
\midrule
October 2019 & 154 & 154 & 0 & 0 & 0 & 0 & 144 & 7 / 3 \\
October 2020 & 182 & 182 & 0 & 0 & 0 & 0 & 167 & 13 / 2 \\
October 2021 & 297 & 297 & 0 & 0 & 0 & 0 & 270 & 16 / 11 \\
October 2022 & 206 & 206 & 0 & 0 & 0 & 0 & 178 & 18 / 10 \\
\midrule
\textbf{Total} & \textbf{839} & \textbf{839} & \textbf{0} & \textbf{0} & \textbf{0} & \textbf{0} & \textbf{759} & \textbf{54 / 26} \\
\bottomrule
\end{tabular}
\end{table}

The prediction file is aligned to the input contract by the stable \texttt{sample\_id}, not by row order. All 839 strict inputs have exactly one matching prediction. There are no missing predictions, unexpected prediction IDs, duplicate IDs, or unknown labels after canonicalisation. This validates the data plumbing, but it does not validate semantic accuracy: without independently annotated gold labels, the prediction distribution cannot be converted into accuracy, macro-F1, MCC, or Cohen's kappa.

\begin{figure}[htbp]
\centering
\includegraphics[width=0.78\textwidth]{images/multiyear_label_distribution.png}
\caption{Composition of the predicted sentiment labels by October period. Neutral predictions dominate every period, while positive and negative groups are comparatively small. The imbalance reduces the precision of group-level comparisons and motivates reporting group sizes together with every estimate.}
\label{fig:multi-year-label-distribution}
\end{figure}

\paragraph{Market data.}
Daily OHLCV data were retrieved through \texttt{vnstock} for the identified securities and the VNINDEX and VN30 reference series. For every October period, the stored price window begins in May of the same year and ends in December. This provides pre-event observations for the market model and post-event observations for lag calculations without importing prices from a different calendar regime.

\begin{table}[htbp]
\centering
\small
\caption{Market-price coverage by October period.}
\label{tab:multi-year-price-coverage}
\begin{tabular}{lrrrrr}
\toprule
\textbf{Period} & \textbf{Price rows} & \textbf{Symbols} & \textbf{Price start} & \textbf{Price end} & \textbf{Strict rows with price} \\
\midrule
October 2019 & 8,320 & 49 & 2019-05-02 & 2019-12-31 & 141 \\
October 2020 & 9,449 & 57 & 2020-05-04 & 2020-12-31 & 168 \\
October 2021 & 11,935 & 73 & 2021-05-04 & 2021-12-31 & 281 \\
October 2022 & 9,182 & 55 & 2022-05-04 & 2022-12-30 & 206 \\
\bottomrule
\end{tabular}
\end{table}

The number of symbols differs by period because the news corpus and entity-linking output differ across years. The price data are therefore comparable in construction, but not identical in cross-sectional composition. This point is important when comparing the magnitude of estimates across periods.

\subsection{Time Alignment and Return Construction}
\label{subsec:time-alignment-multi-year}

The data contain daily prices but not the exact time at which an article was disseminated to market participants. The implemented alignment rule maps each article's normalized publication date to the first trading date strictly after that calendar date. Thus, an article published on a trading day is evaluated against the next available trading session; an article published on a weekend or holiday is also mapped to the next trading session. No fixed 15:00 cutoff is imposed, because the daily data do not identify the intraday time at which the information was incorporated into the price.

For security $i$ on trading date $t$, the continuously compounded percentage return is computed as
\begin{equation}
 r_{i,t}=100\times\left[\log(P_{i,t})-\log(P_{i,t-1})\right],
\end{equation}
where $P_{i,t}$ is the recorded closing price. The market return is computed using the same transformation for VNINDEX. The analysis uses the downloaded price series as stored; any interpretation of return magnitude is conditional on the source's adjustment and corporate-action coverage.

The sentiment score is a deterministic label encoding:
\begin{equation}
 S_{i,t}=\begin{cases}
 +1,&\text{positive prediction},\\
 0,&\text{neutral prediction},\\
 -1,&\text{negative prediction}.
 \end{cases}
\end{equation}

Several article--ticker rows can refer to the same ticker on the same date. Because the market supplies one return observation per ticker--day, the financial unit is formed by aggregation before inference. The panel records mean sentiment, positive/neutral/negative shares, the number of unique articles, and the number of contributing article--ticker rows. This avoids artificially multiplying the effective sample size by repeating one return for every article.

\subsection{Event-Study Methodology}
\label{subsec:event-study-multi-year}

Following the standard market-model event-study framework of Brown and Warner~\cite{brownwarner1985}, the expected return for each security is estimated from observations before 1 October of the relevant year:
\begin{equation}
 R_{i,t}=\alpha_i+\beta_iR_{m,t}+\varepsilon_{i,t},
\end{equation}
where $R_{m,t}$ is the VNINDEX return. The abnormal return is
\begin{equation}
 AR_{i,t}=R_{i,t}-\left(\widehat{\alpha}_i+\widehat{\beta}_iR_{m,t}\right).
\end{equation}
A security-specific market model is fitted when at least 20 pre-event observations are available; otherwise the documented fallback parameters are retained. The final period-specific runs contain 49, 57, 73, and 55 stored price symbols, respectively, with 54 fitted market-model records in the October 2022 run and corresponding period-specific records in the other runs.

For an event aligned to trading date $t$, cumulative abnormal return is defined as
\begin{equation}
 CAR_i(T_1,T_2)=\sum_{\tau=T_1}^{T_2}AR_{i,t+\tau}.
\end{equation}
Three windows are examined: $[0,0]$ for the event session, $[0,+1]$ for the event session and the next session, and $[0,+3]$ for the event session and the next three sessions. For every period and window, the pipeline reports the mean and median CAR, a two-sided one-sample $t$-test against zero, a two-sided sign-flip permutation test with 5,000 replications, and a percentile bootstrap 95\% confidence interval with 5,000 replications. Fixed random seeds make the resampling results reproducible.

\begin{figure}[htbp]
\centering
\includegraphics[width=0.84\textwidth]{images/multiyear_car_by_period.png}
\caption{Aggregate CAR by October period and event window. Points show mean CAR and whiskers show percentile bootstrap 95\% confidence intervals. The figure is intended for cross-period comparison; it does not apply a multiple-comparison correction and does not establish causality.}
\label{fig:multi-year-car}
\end{figure}

\begin{longtable}{llrrrrrr}
\caption{Aggregate event-study CAR tests by period and event window.}\label{tab:multi-year-car-tests}\\
\toprule
\textbf{Period} & \textbf{Window} & \textbf{N} & \textbf{Mean} & \textbf{Median} & \textbf{$p_t$} & \textbf{$p_{perm}$} & \textbf{Bootstrap 95\% CI} \\
\midrule
\endfirsthead
\toprule
\textbf{Period} & \textbf{Window} & \textbf{N} & \textbf{Mean} & \textbf{Median} & \textbf{$p_t$} & \textbf{$p_{perm}$} & \textbf{Bootstrap 95\% CI} \\
\midrule
\endhead
\midrule
\multicolumn{8}{r}{Continued on next page}\\
\endfoot
\bottomrule
\endlastfoot
2019-10 & $[0,0]$ & 120 & -0.1323 & -0.1570 & 0.5420 & 0.5491 & $[-0.5550,0.2916]$ \\
2019-10 & $[0,+1]$ & 120 & 0.0573 & -0.1285 & 0.8833 & 0.8830 & $[-0.6841,0.8108]$ \\
2019-10 & $[0,+3]$ & 120 & 0.0020 & -0.2812 & 0.9970 & 0.9974 & $[-1.0491,1.0923]$ \\
2020-10 & $[0,0]$ & 134 & 0.0522 & -0.2528 & 0.8031 & 0.8072 & $[-0.3645,0.4595]$ \\
2020-10 & $[0,+1]$ & 134 & -0.2975 & -0.6585 & 0.3424 & 0.3363 & $[-0.8991,0.3050]$ \\
2020-10 & $[0,+3]$ & 133 & -0.9261 & -1.3113 & 0.0438 & 0.0440 & $[-1.8254,-0.0657]$ \\
2021-10 & $[0,0]$ & 212 & -0.2797 & -0.4321 & 0.1070 & 0.1096 & $[-0.6117,0.0587]$ \\
2021-10 & $[0,+1]$ & 212 & -0.2584 & -0.5174 & 0.2920 & 0.2953 & $[-0.7377,0.2377]$ \\
2021-10 & $[0,+3]$ & 212 & -1.0149 & -1.0412 & 0.0101 & 0.0088 & $[-1.7967,-0.2662]$ \\
2022-10 & $[0,0]$ & 150 & 0.1363 & 0.0008 & 0.6023 & 0.6091 & $[-0.3557,0.6421]$ \\
2022-10 & $[0,+1]$ & 150 & -0.0220 & 0.1798 & 0.9501 & 0.9562 & $[-0.7298,0.6360]$ \\
2022-10 & $[0,+3]$ & 150 & -0.2889 & -0.7977 & 0.5614 & 0.5615 & $[-1.2725,0.6793]$ \\
\end{longtable}

\paragraph{Event-study results.}
The aggregate CAR results do not support a single stable sentiment-linked price response across the four October periods. The same-day window is statistically indistinguishable from zero in every period: all $p$-values exceed 0.10 and all bootstrap intervals include zero. In the three-session window, 2020-10 and 2021-10 produce negative mean CAR estimates with nominal $p$-values below 0.05. These two results are noteworthy but should not be treated as decisive evidence: they are two findings among many period--window tests, they are not adjusted for multiple comparisons, and the negative three-session CAR is not itself evidence that sentiment caused the decline.

The 2019 and 2022 estimates remain close to zero or statistically imprecise across all windows. This cross-period heterogeneity is more consistent with a period-dependent market association than with a universal FinABSA effect. It may reflect differences in market regime, news composition, entity coverage, macroeconomic shocks, or the composition of the available price universe.

\subsection{News-Day Versus No-News-Day Control}
\label{subsec:news-control-multi-year}

The event study is complemented by a control comparison. For each period, the control sample contains eligible ticker--days belonging to tickers that appear in the event sample, restricted to October of that year. A news ticker--day is a ticker--day with at least one aligned article event; a no-news ticker--day is an eligible ticker--day without an identified event. The estimand is
\begin{equation}
 \Delta=\overline{AR}_{news}-\overline{AR}_{no\text{-}news}.
\end{equation}

The two distributions are compared using a Welch unequal-variance test, a pooled-label permutation test with 5,000 permutations, and a percentile bootstrap interval for the difference in means.

\begin{table}[htbp]
\centering
\scriptsize
\caption{News-day versus eligible no-news-day abnormal returns.}
\label{tab:news-control-multi-year}
\begin{tabular}{lrrrrrrrr}
\toprule
\textbf{Period} & \textbf{$N_{news}$} & \textbf{$N_{0}$} & \textbf{News mean} & \textbf{No-news mean} & \textbf{$\Delta$} & \textbf{Welch $p$} & \textbf{Perm. $p$} & \textbf{Bootstrap 95\% CI} \\
\midrule
2019-10 & 120 & 941 & -0.1323 & 0.0064 & -0.1386 & 0.5383 & 0.4585 & $[-0.5937,0.3095]$ \\
2020-10 & 134 & 1,027 & 0.0522 & -0.2161 & 0.2683 & 0.2240 & 0.1828 & $[-0.1601,0.7062]$ \\
2021-10 & 212 & 1,217 & -0.2797 & -0.1039 & -0.1758 & 0.3553 & 0.3771 & $[-0.5493,0.1980]$ \\
2022-10 & 150 & 936 & 0.1363 & -0.3559 & 0.4921 & 0.0804 & 0.0684 & $[-0.0645,1.0442]$ \\
\bottomrule
\end{tabular}
\end{table}

\begin{figure}[htbp]
\centering
\includegraphics[width=0.78\textwidth]{images/multiyear_news_day_control.png}
\caption{News-day minus eligible no-news-day abnormal return by October period. Whiskers show percentile bootstrap 95\% confidence intervals; each interval includes zero.}
\label{fig:multi-year-news-control}
\end{figure}

The control comparison does not reject equality at the 5\% level in any period. The largest positive point estimate occurs in October 2022, where news days have 0.4921 percentage points higher abnormal returns than no-news days, but the Welch and permutation $p$-values are 0.0804 and 0.0684, respectively, and the bootstrap interval crosses zero. The result is therefore suggestive but not conventionally significant. The comparison remains observational because news selection is not random: firms that receive coverage may systematically differ from firms that do not.

\subsection{Sentiment-Group Contrasts}
\label{subsec:sentiment-groups-multi-year}

At the ticker--day level, the analysis compares three pre-specified contrasts: positive minus neutral, negative minus neutral, and positive minus negative. Each contrast is evaluated with a Welch test, a pooled permutation test, and a bootstrap interval. Differences are defined as group A minus group B.

\begin{longtable}{llrrrrrr}
\caption{Sentiment-group abnormal-return contrasts by period.}\label{tab:sentiment-contrasts-multi-year}\\
\toprule
\textbf{Period} & \textbf{Contrast} & \textbf{$N_A$} & \textbf{$N_B$} & \textbf{Difference} & \textbf{Welch $p$} & \textbf{Perm. $p$} & \textbf{Bootstrap 95\% CI} \\
\midrule
\endfirsthead
\toprule
\textbf{Period} & \textbf{Contrast} & \textbf{$N_A$} & \textbf{$N_B$} & \textbf{Difference} & \textbf{Welch $p$} & \textbf{Perm. $p$} & \textbf{Bootstrap 95\% CI} \\
\midrule
\endhead
\midrule
\multicolumn{8}{r}{Continued on next page}\\
\endfoot
\bottomrule
\endlastfoot
2019-10 & Positive $-$ Neutral & 6 & 113 & -1.5118 & 0.1196 & 0.1252 & $[-3.2365,-0.3238]$ \\
2019-10 & Negative $-$ Neutral & 1 & 113 & -- & -- & -- & Not estimable \\
2019-10 & Positive $-$ Negative & 6 & 1 & -- & -- & -- & Not estimable \\
2020-10 & Positive $-$ Neutral & 10 & 123 & -0.8544 & 0.3380 & 0.2813 & $[-2.1566,0.9352]$ \\
2020-10 & Negative $-$ Neutral & 1 & 123 & -- & -- & -- & Not estimable \\
2020-10 & Positive $-$ Negative & 10 & 1 & -- & -- & -- & Not estimable \\
2021-10 & Positive $-$ Neutral & 12 & 191 & 1.5693 & 0.0267 & 0.0388 & $[0.5210,2.8516]$ \\
2021-10 & Negative $-$ Neutral & 9 & 191 & 0.0915 & 0.9093 & 0.9112 & $[-1.0946,1.7607]$ \\
2021-10 & Positive $-$ Negative & 12 & 9 & 1.4779 & 0.1460 & 0.1372 & $[-0.4072,3.2104]$ \\
2022-10 & Positive $-$ Neutral & 14 & 126 & 0.6990 & 0.4518 & 0.4271 & $[-0.9308,2.4970]$ \\
2022-10 & Negative $-$ Neutral & 10 & 126 & -0.4688 & 0.6712 & 0.6601 & $[-2.5243,1.4876]$ \\
2022-10 & Positive $-$ Negative & 14 & 10 & 1.1678 & 0.3962 & 0.3971 & $[-1.2938,3.6937]$ \\
\end{longtable}

The most notable group contrast appears in October 2021: the positive group has a 1.5693 percentage-point higher mean abnormal return than the neutral group, with Welch $p=0.0267$, permutation $p=0.0388$, and a bootstrap interval that excludes zero. This result should be described as a period-specific exploratory finding rather than a general law. The positive group contains only 12 ticker--day observations, the negative group contains nine, the analysis involves multiple comparisons, and the effect is not replicated in the corresponding positive--neutral comparisons for 2019, 2020, or 2022.

The 2019 positive--neutral point estimate is negative, whereas the 2020 and 2022 differences are also statistically imprecise. The failure of sign consistency across years weakens the claim that predicted sentiment produces a stable cross-sectional ranking of abnormal returns.

\begin{figure}[htbp]
\centering
\includegraphics[width=0.78\textwidth]{images/multiyear_positive_neutral_contrast.png}
\caption{Positive-minus-neutral abnormal-return contrast by October period. Whiskers show bootstrap 95\% confidence intervals. The positive 2021 estimate is not sufficient on its own to establish a stable sentiment effect across years.}
\label{fig:multi-year-positive-neutral}
\end{figure}

\subsection{Two-Way Fixed-Effects Panel Regression}
\label{subsec:panel-regression-multi-year}

To control for persistent differences between securities and common shocks across dates, the following two-way fixed-effects model is estimated separately for each October period:
\begin{equation}
 AR_{i,\tau(i,t)}=\alpha_i+\delta_{\tau(i,t)}+\beta_1S_{i,t}
 +\beta_2\mathrm{NegativeShare}_{i,t}
 +\beta_3\mathrm{PositiveShare}_{i,t}
 +\gamma\log(1+\mathrm{ArticleCount}_{i,t})+\varepsilon_{i,t},
\end{equation}
where $\alpha_i$ denotes ticker fixed effects, $\delta_{\tau(i,t)}$ denotes date fixed effects, and $\tau(i,t)$ is the first trading session strictly after publication. The dependent variable is market-model abnormal return on the aligned evaluation date.

\begin{table}[htbp]
\centering
\scriptsize
\caption{Key coefficients from period-specific two-way fixed-effects regressions. The full coefficient tables, including ticker and date fixed-effect indicators, are provided in the accompanying CSV output.}
\label{tab:panel-key-coefficients-multi-year}
\begin{tabular}{llrrrr}
\toprule
\textbf{Period} & \textbf{Variable} & \textbf{Coefficient} & \textbf{Std. error} & \textbf{$t$-statistic} & \textbf{$p$-value} \\
\midrule
2019-10 & Sentiment score & -0.4450 & 0.8508 & -0.5230 & 0.6009 \\
2019-10 & Negative share & 1.5084 & 1.8412 & 0.8192 & 0.4127 \\
2019-10 & Positive share & 1.0633 & 1.2562 & 0.8468 & 0.3973 \\
2019-10 & Log article count & -1.5405 & 3.3930 & -0.4540 & 0.6498 \\
2020-10 & Sentiment score & 2.5056 & 0.8042 & 3.1156 & 0.0018 \\
2020-10 & Negative share & -4.9332 & 1.8033 & -2.7356 & 0.0062 \\
2020-10 & Positive share & -2.4276 & 1.2770 & -1.9010 & 0.0573 \\
2020-10 & Log article count & 2.3897 & 1.9986 & 1.1957 & 0.2318 \\
2021-10 & Sentiment score & -0.9696 & 2.4681 & -0.3927 & 0.6944 \\
2021-10 & Negative share & 0.1799 & 2.5174 & 0.0715 & 0.9430 \\
2021-10 & Positive share & 0.4345 & 2.6128 & 0.1663 & 0.8679 \\
2021-10 & Log article count & 0.5867 & 1.0124 & 0.5795 & 0.5623 \\
2022-10 & Sentiment score & 6.5062 & 5.3763 & 1.2100 & 0.2262 \\
2022-10 & Negative share & 5.9827 & 5.2919 & 1.1305 & 0.2582 \\
2022-10 & Positive share & -6.1476 & 5.5086 & -1.1160 & 0.2644 \\
2022-10 & Log article count & 0.3079 & 1.3849 & 0.2224 & 0.8241 \\
\bottomrule
\end{tabular}
\end{table}

The sentiment-score coefficient is positive and statistically significant only in the 2020 period. The coefficient is negative and insignificant in 2019 and 2021, and positive but imprecise in 2022. This sign instability is central: a robust sentiment--return association would be expected to exhibit a more consistent direction under comparable definitions. The 2020 coefficient therefore constitutes a period-specific association, not evidence that the model has a stable return-prediction relationship.

The coefficient on negative share is also significant only in 2020 and is negative, while the positive-share coefficient is not significant at the 5\% level. Because sentiment shares are compositional and mechanically related to the omitted neutral share, their coefficients should not be interpreted as independent causal marginal effects. The relatively high within-period $R^2$ values are not evidence of strong predictive power: ticker and date fixed effects can absorb substantial variation, while the sentiment covariates remain sensitive to sample composition.

The intended covariance estimator was HC3. In sparse two-way fixed-effects designs, HC3 can produce infinite standard errors when leverage reaches a singular boundary. The implementation therefore falls back to HC1 when necessary and records this fallback in the period-specific metadata. These panel coefficients should consequently be treated as exploratory conditional associations.

\begin{figure}[htbp]
\centering
\includegraphics[width=0.78\textwidth]{images/multiyear_panel_sentiment_coefficient.png}
\caption{Period-specific sentiment-score coefficients from the two-way fixed-effects models. Whiskers show coefficient plus or minus 1.96 standard errors. The 2020 coefficient is positive and relatively precise, but the sign is not replicated across the other periods.}
\label{fig:multi-year-panel-coefficient}
\end{figure}

\subsection{Lagged Robustness and Sensitivity Analysis}
\label{subsec:lag-robustness-multi-year}

Lag analysis evaluates abnormal returns at 0, 1, 2, 3, 5, and 10 trading sessions after the aligned event date. These are sensitivity checks rather than independent causal tests. The sample size decreases at longer lags because the relevant security price must remain available after the event.

\begin{longtable}{llrrrrrr}
\caption{Lagged abnormal-return robustness by period.}\label{tab:lag-robustness-multi-year}\\
\toprule
\textbf{Period} & \textbf{Lag} & \textbf{N} & \textbf{Overall AR} & \textbf{Negative} & \textbf{Neutral} & \textbf{Positive} \\
\midrule
\endfirsthead
\toprule
\textbf{Period} & \textbf{Lag} & \textbf{N} & \textbf{Overall AR} & \textbf{Negative} & \textbf{Neutral} & \textbf{Positive} \\
\midrule
\endhead
\midrule
\multicolumn{7}{r}{Continued on next page}\\
\endfoot
\bottomrule
\endlastfoot
2019-10 & 0 & 120 & -0.1323 & -1.5766 & -0.0648 & 0.8940 \\
2019-10 & 1 & 146 & 0.3194 & 0.0588 & 0.2772 & 1.2785 \\
2019-10 & 2 & 141 & -0.0968 & 0.0331 & -0.0974 & -0.1286 \\
2019-10 & 3 & 134 & 0.0067 & -0.4507 & 0.0283 & -0.3248 \\
2019-10 & 5 & 108 & 0.1154 & -0.6470 & 0.1166 & 0.2232 \\
2019-10 & 10 & 74 & -0.1635 & 0.1766 & -0.1547 & -0.6273 \\
2020-10 & 0 & 134 & 0.2261 & -1.9519 & 0.3085 & -0.4498 \\
2020-10 & 1 & 164 & -0.1520 & 0.5470 & -0.1975 & 0.2915 \\
2020-10 & 2 & 159 & -0.3772 & 1.3165 & -0.3780 & -0.7051 \\
2020-10 & 3 & 152 & 0.0687 & 0.9315 & 0.1112 & -0.6345 \\
2020-10 & 5 & 144 & -0.0997 & 0.4069 & -0.1472 & 0.4869 \\
2020-10 & 10 & 100 & -0.4924 & 0.3574 & -0.5596 & 1.5471 \\
2021-10 & 0 & 212 & -0.0285 & -0.1815 & -0.0939 & 1.0708 \\
2021-10 & 1 & 267 & 0.2065 & 0.6397 & 0.2368 & -0.5095 \\
2021-10 & 2 & 255 & -0.3156 & -0.4267 & -0.2764 & -0.8200 \\
2021-10 & 3 & 244 & 0.0773 & -0.3897 & 0.1940 & -1.2377 \\
2021-10 & 5 & 219 & 0.1398 & -0.1027 & 0.1966 & -0.4387 \\
2021-10 & 10 & 147 & -0.0471 & 1.2318 & -0.1101 & -0.1871 \\
2022-10 & 0 & 150 & -0.0936 & -0.7884 & -0.1653 & 0.8940 \\
2022-10 & 1 & 174 & -0.1797 & 0.3368 & 0.0713 & -2.7331 \\
2022-10 & 2 & 163 & -0.2007 & -1.2344 & -0.0991 & -0.4067 \\
2022-10 & 3 & 153 & -0.5972 & -2.6554 & -0.4837 & -0.1760 \\
2022-10 & 5 & 134 & -0.7243 & 0.4479 & -0.9416 & 0.5251 \\
2022-10 & 10 & 93 & -0.5183 & -0.4422 & -0.2779 & -4.1530 \\
\end{longtable}

\begin{figure}[htbp]
\centering
\includegraphics[width=0.84\textwidth]{images/multiyear_lag_robustness.png}
\caption{Lagged sensitivity of mean abnormal returns across the four October periods. The changing signs and magnitudes indicate that the association is sensitive to the evaluation horizon rather than a stable delayed response.}
\label{fig:multi-year-lag-robustness}
\end{figure}

The lag results do not display a common monotonic decay. In 2019, the overall estimate changes from negative at lag 0 to positive at lag 1 and back toward zero thereafter. In 2020, the estimates become more negative at lags 2 and 10, while in 2021 the sign alternates across horizons. In 2022, the overall mean is increasingly negative at lags 3 and 5. Sentiment-specific estimates are even less stable, particularly for the small positive and negative groups. The appropriate conclusion is that the apparent price association depends on the timing convention and the available evaluation horizon.

\subsection{Additional Diagnostic Evidence}
\label{subsec:additional-diagnostics-multi-year}

The output package also contains daily ticker--day sentiment and abnormal-return tables, market-model coverage, event observations, and panel-regression data. These artifacts support auditability and allow the reader to trace each reported result back to an aligned ticker--day observation. The following diagnostic figures are descriptive rather than inferential.

\begin{figure}[htbp]
\centering
\includegraphics[width=0.82\textwidth]{images/market_model_coverage.png}
\caption{Market-model coverage diagnostic from the completed October 2022 run. It shows the number of pre-event observations available for individual symbols and illustrates why coverage thresholds and fallback handling are necessary.}
\label{fig:market-model-coverage-multi-year}
\end{figure}

\begin{figure}[htbp]
\centering
\includegraphics[width=0.90\textwidth]{images/sentiment_vs_abnormal_return.png}
\caption{Descriptive comparison between daily mean sentiment and mean abnormal return in the ticker--day panel. The figure is not a causal identification plot and does not control for all market-wide shocks.}
\label{fig:sentiment-vs-ar-multi-year}
\end{figure}

\subsection{Multiple Comparisons and Statistical Interpretation}
\label{subsec:multiple-comparisons}

The multi-year design produces several families of related tests: four periods, three event windows, multiple sentiment groups, a news-day control, panel coefficients, and six lag horizons. Consequently, at least some nominally small $p$-values can occur by chance when many hypotheses are examined. The results are therefore interpreted using three criteria: the size and direction of the estimate, whether the confidence interval excludes zero, and whether the finding is replicated under related specifications or periods.

Under this framework, the 2020 and 2021 three-day CAR findings and the 2021 positive--neutral contrast are interesting signals for further study, but not definitive discoveries. They should be reported with their exact sample sizes, intervals, and unadjusted $p$-values. A confirmatory study should pre-specify a smaller number of primary hypotheses and apply a family-wise or false-discovery-rate correction, such as Holm or Benjamini--Hochberg, depending on the inferential objective.

\subsection{Analyses Not Validly Estimated in the Current Run}
\label{subsec:deferred-analyses-multi-year}

Several analyses appeared in the initial experimental plan but are not assigned numerical claims here because the required artifacts are absent. First, intrinsic CafeF accuracy, macro-F1, MCC, and Cohen's kappa require an independently human-annotated gold-label sample. The model's own predictions cannot serve as their ground truth. Second, probability calibration metrics such as Expected Calibration Error require valid per-class probabilities or logits; the current prediction files contain hard labels and empty probability fields. Third, out-of-time return forecasting requires a pre-specified chronological training, validation, and test design with a fitted return-prediction model; the current pipeline is an event-study association analysis, not a trained forecaster. Fourth, a broad-linking robustness result requires a complete independent inference and financial analysis run on the 2,740-row broad input set; strict and full input files alone are not sufficient.

This separation is methodologically important. It prevents unsupported accuracy, calibration, forecasting, and robustness numbers from being presented as results. The completed outputs can be extended later without changing the core data contract.

\subsection{Limitations}
\label{subsec:limitations-statistical}

\paragraph{Limited temporal coverage and regime dependence.}
Although the study now covers four October periods, it still samples only one month per year. October 2019--2022 may contain unusual news intensity, market stress, or sector composition. The estimates should not be interpreted as representative of every month or every COVID-19 phase. The common-month design improves comparability but does not eliminate regime dependence.

\paragraph{Single-source publication selection.}
CafeF is one news source, and its editorial selection is not random. The presence, timing, and wording of an article may depend on firm size, liquidity, market attention, or an ongoing price movement. The news-day control cannot fully remove this selection process.

\paragraph{Entity-linking measurement error.}
The strict threshold improves precision but reduces coverage. The broad candidate file may contain false positives from uppercase abbreviations, while strict filtering may exclude valid ticker mentions. A linking error can connect sentiment to the wrong price series or remove a relevant event entirely. The 32 failed price symbols and 43 strict rows without usable price windows make this limitation observable rather than hypothetical.

\paragraph{Class imbalance and model uncertainty.}
Across the 839 strict predictions, 759 are neutral, 54 are positive, and 26 are negative. Small positive and negative groups produce wide intervals and make group comparisons sensitive to individual observations. Moreover, the current outputs are hard labels without calibrated probabilities, so uncertainty in the classifier is not propagated into the financial tests.

\paragraph{Daily timing and return measurement.}
Daily prices do not reveal whether an article was published before or after the market had already incorporated the information. Mapping every article to the next trading session is conservative and reproducible, but it may delay some events or miss same-day reactions. The price-source adjustment and corporate-action history are also not independently audited in this dataset.

\paragraph{Sparse panel and covariance estimation.}
The period-specific panels have different numbers of ticker--day observations and dates. Sparse fixed-effects designs can produce high leverage and unstable standard errors. The HC1 fallback after a singular HC3 correction is documented, but it reduces the strength of any regression-based inference. The coefficients should therefore be treated as exploratory conditional associations.

\paragraph{Non-causal design and multiple testing.}
The analysis is observational and examines many specifications. It does not control for every confounder, does not use an exogenous news shock, and does not pre-register all hypotheses. Statistical significance in one period or window is not sufficient to establish a causal effect or a deployable trading signal.

\subsection{Future Work}
\label{subsec:future-work-statistical}

Future work should proceed in a pre-specified sequence. First, create an independently annotated and adjudicated CafeF gold set with multiple annotators, report inter-annotator agreement, and evaluate accuracy, macro-F1, balanced accuracy, MCC, and class-specific confusion patterns. The gold set should be split by article or time period so that repeated articles do not leak across train and test data.

Second, repeat inference for all four periods using a frozen model checkpoint and preserve per-class probabilities or logits. This will enable confidence-weighted aggregation, Expected Calibration Error, reliability diagrams, and sensitivity analysis that separates high-confidence from low-confidence predictions.

Third, expand the time axis beyond one month per year, ideally using multiple months before, during, and after the pandemic regime. A longer panel would increase power, allow year and sector fixed effects, and support formal interactions between sentiment and market regime. The broad-linking run should be completed independently and compared with the strict run using pre-defined matching rules.

Fourth, strengthen the financial design by estimating alternative market-adjustment models, including market-adjusted returns and factor models, while fixing the estimation-window rule before inspecting results. Placebo dates, non-news dates, sector decomposition, and clustered standard errors should be included as confirmatory robustness checks rather than selected after seeing the main results.

Finally, if return prediction is a research objective, train a separate chronological forecasting model with a censor gap between the information date and evaluation window. Performance should be compared with simple baselines using out-of-time RMSE, MAE, rank correlation, and economically interpretable evaluation metrics. Any trading simulation would require transaction costs, liquidity constraints, delisting treatment, and strict protection against look-ahead bias.

\subsection{Overall Assessment and Conclusion}
\label{subsec:conclusion-statistical}

The completed analysis establishes a reproducible end-to-end connection between Vietnamese financial headlines and market outcomes. It validates the prediction contract for 839 strict article--ticker inputs across October 2019--2022, constructs target-specific sentiment signals, aligns publication dates to the next trading session, estimates security-level market models, and applies complementary event-study, control-group, group-comparison, panel, and lagged procedures.

The financial evidence is mixed and period-dependent. Aggregate same-day CAR is not statistically different from zero in any year. Negative three-session CAR appears in 2020 and 2021 at nominal significance levels, while the 2021 positive--neutral contrast and the 2020 panel sentiment coefficient are also nominally significant. However, the signs and magnitudes are not stable across years, event windows, or lags; the significant findings arise within a large family of tests; and the confidence intervals and small sentiment groups indicate substantial uncertainty. The news-day control does not reach the conventional 5\% threshold in any period.

The most defensible conclusion is therefore that the pipeline is technically complete and produces observable sentiment--return associations, but the current evidence does not demonstrate a stable, causal, or reliably predictive effect of FinABSA sentiment on Vietnamese stock returns. The occasional significant estimates are useful hypotheses for future confirmatory research, not evidence of a standalone investment strategy.

\begin{table}[htbp]
\centering
\small
\caption{Reproducibility checklist for the completed multi-year validation.}
\label{tab:reproducibility-multi-year}
\begin{tabular}{p{0.42\textwidth}p{0.46\textwidth}}
\toprule
\textbf{Component} & \textbf{Completed implementation} \\
\midrule
Sampling design & October 2019, October 2020, October 2021, October 2022 \\
Price window & May--December of the same year for every period \\
Prediction alignment & Stable \texttt{sample\_id}; 839/839 matched \\
Label normalization & Canonical positive/neutral/negative; zero unknown labels \\
Entity-linking primary sample & Strict confidence threshold $\geq 0.95$ \\
Event mapping & First trading session strictly after publication date \\
Market adjustment & Security-specific market model with VNINDEX proxy \\
Event windows & $[0,0]$, $[0,+1]$, and $[0,+3]$ \\
Inference tests & One-sample $t$-test, sign-flip permutation, bootstrap CI \\
Control test & News-day versus eligible no-news-day comparison \\
Group tests & Welch, permutation, and bootstrap contrasts \\
Panel model & Ticker/date fixed effects; HC1 fallback documented \\
Lag sensitivity & 0, 1, 2, 3, 5, and 10 trading-session lags \\
Output artifacts & CSV tables, period-specific event files, six multi-year figures, Markdown report, Excel workbook \\
\bottomrule
\end{tabular}
\end{table}

The numerical source files are stored under \texttt{data/cafef\_multiyear/analysis\_multiyear/}. The principal tables are \texttt{car\_primary\_all.csv}, \texttt{news\_day\_control\_tests.csv}, \texttt{sentiment\_group\_tests.csv}, \texttt{panel\_key\_coefficients.csv}, \texttt{robustness\_summary.csv}, and \texttt{prediction\_validation\_clean.csv}. The complete workbook is \texttt{multiyear\_validation\_tables.xlsx}, and the period-specific diagnostics are retained under each period's \texttt{analysis/} directory.

\begin{quote}
This is research and analysis only, not personalized financial advice.
\end{quote}
